\documentclass[11pt]{article}
\usepackage[margin=1in]{geometry}
\usepackage{amsmath,amssymb,amsthm,mathtools,bm,booktabs,enumitem,microtype}
\usepackage[hidelinks]{hyperref}

\newtheorem{theorem}{Theorem}[section]
\newtheorem{proposition}[theorem]{Proposition}
\newtheorem{corollary}[theorem]{Corollary}
\newtheorem{lemma}[theorem]{Lemma}
\theoremstyle{definition}
\newtheorem{definition}[theorem]{Definition}
\newtheorem{remark}[theorem]{Remark}

\newcommand{\Hh}{\mathcal H}
\newcommand{\Th}{\Theta}
\newcommand{\R}{\mathcal R}
\newcommand{\ip}[2]{\langle #1,#2\rangle}
\newcommand{\norm}[1]{\left\lVert #1\right\rVert}
\newcommand{\op}{\mathrm{op}}
\newcommand{\Ran}{\mathrm{Ran}}
\newcommand{\spec}{\mathrm{spec}}
\newcommand{\Tr}{\operatorname{Tr}}
\newcommand{\Id}{I}

\title{When Is a Neural Network Ready to Become Smaller?\\
Safe Dynamic Spectral Reduction from Target-Weighted Prediction Geometry}
\author{Enrico Bignozzi}
\date{Working theory draft --- September 2026}

\begin{document}
\maketitle

\begin{abstract}
Overparameterization can be useful early in training because it lets a network discover a representation, yet the same excess capacity can become statistically and computationally wasteful after target-relevant directions have separated from nuisance directions. This note develops a theory for deciding when capacity may be removed during training. The primitive object is the prediction Jacobian, not the spectrum of the weight matrices. Under square loss, the gradient energy carried by a prediction mode is exactly $\mu_j a_j^2$, where $\mu_j$ is a prediction-operator eigenvalue and $a_j$ is residual target mass in its eigenvector. In fixed geometry, the exact finite-horizon value of keeping that mode is $\frac12 a_j^2(1-e^{-2\mu_jH})$. This yields an exact value-per-compute oracle and shows why small eigenvalues alone are not a pruning criterion. A compute-to-accuracy theorem makes the optimization claim precise: truncation can reduce both per-step cost and the scalar-step condition-number bound, but only for target accuracies above the deletion floor. An exact break-even identity then separates the same-step discovery advantage created by temporary width from the extra optimization progress a compact-from-start model can buy with the saved compute. Thus safe contraction is necessary but not sufficient for a learn-wide-then-shrink compute advantage. Prospectively, an exact local theorem shows that a safe reduced update is more efficient for the next infinitesimal unit of compute exactly when the fraction of target-weighted gradient energy it discards is smaller than the fraction of per-update compute it saves. In frozen geometry this extends to every finite compute budget: contraction wins exactly when the extra progress bought on retained modes exceeds the progress forfeited on deleted modes.

For moving nonlinear geometry we compare the full gradient flow with a projected flow and obtain an explicit finite-horizon trajectory and risk bound in terms of omitted gradient energy. Curvature and eigengap arguments control the future growth and rotation of a discarded spectral cluster. A structural-transfer theorem then separates the ideal spectral decision from its realization by neurons, channels, heads, or low-rank factors, and a master theorem combines these ingredients into one finite-horizon risk certificate. A checkpoint-computable refinement removes the unknown future loss drop and turns spectral-gap persistence itself into a current-state gate. For the central lower-bulk use case, a target-mass persistence theorem further separates three objects that eigenvalue thresholding conflates: the bulk edge, residual target mass still trapped in the bulk, and the amount of spectral motion that can move new target mass into it. Their combination yields an explicit geometric bulk-onset time. An additive rank-factor parameterization then supplies an exact physical structural bridge: deleting a rank-one factor has an exact ghost embedding, an explicit contraction jump, and an exact group-gradient-energy budget. This yields a complementary gap-free safety certificate for literal parameter groups. These population prediction-gap bounds feed a geometry-localized finite-sample damage certificate. For the final hidden layer with a linear readout, a rank-$r$ representation always contains a subset of $r$ neurons spanning the same current readout space; the only reason not to contract immediately is future feature-discovery value.

Finally, the exact quadratic teacher--student phase theorem gives a nonzero endogenous onset: target spikes are initially hidden in a random bulk, become identifiable only after dynamic BBP crossings, and eventually span an asymptotically target-sufficient $k$-dimensional prediction subspace while the nuisance bulk contracts toward zero. In the rank-one model, contracting blindly before separation induces a $\log d$ target-discovery delay, whereas retaining the post-BBP outlier has order-one alignment. The rank-one radial dynamics are also exactly solvable, producing a critical wide-to-compact compute ratio $\rho_\star$: temporary width repays its discovery phase precisely when the realized structural cost ratio lies below this frontier. The theory therefore predicts a two-stage training principle: overparameterize to discover, then contract after target-relevant geometry becomes resolved, structurally realizable, target-compressible, statistically observable, and cheap to discard in finite-horizon value. The break-even condition adds the final requirement for a compute claim: the discovery benefit generated by temporary width must exceed its incremental compute opportunity cost.
\end{abstract}

\section{Problem and principle}
A dense model contains many parameter directions that eventually exert almost no useful effect on prediction. The tempting rule
\[
\mu_j(t)\approx 0 \quad\Longrightarrow\quad \text{delete mode }j
\]
is not safe. A weak mode may carry important unresolved target signal, and in a feature-learning network a direction that is currently buried in a random bulk can later become a target-aligned outlier. Conversely, a large mode can be pure nuisance capacity.

The goal is therefore not generic pruning. It is to answer three separate questions:
\begin{enumerate}[label=(\roman*)]
\item \emph{Current value:} how much terminal risk reduction can a mode buy if the geometry is frozen now?
\item \emph{Future discovery value:} how much could the geometry associated with that mode improve over the next training horizon?
\item \emph{Structural realizability:} can the low-value spectral sector be removed as actual neurons, channels, heads, blocks, or rank without destroying the kept sector?
\end{enumerate}
The working principle is
\[
\boxed{\text{overparameterize to discover; contract only after useful geometry has separated.}}
\]

\section{Prediction-space geometry}
Let $F:\Th\to\Hh$ be a differentiable predictor into a real Hilbert prediction space. For square loss define
\[
e_\theta=f^\star-F(\theta),\qquad \R(\theta)=\frac12\norm{e_\theta}^2.
\]
Write $J_\theta=DF(\theta):\Th\to\Hh$. Population gradient flow is
\[
\dot\theta_t=J_t^\ast e_t,
\qquad
\dot e_t=-L_t e_t,
\qquad
L_t=J_tJ_t^\ast.
\]
The weight spectrum, the parameter Gram spectrum $J_t^\ast J_t$, and the prediction spectrum $L_t$ are distinct objects. The nonzero spectra of $J_t^\ast J_t$ and $L_t$ coincide, but weight singular values do not in general determine either.

Assume for the moment that $L_t$ has simple positive eigenpairs
\[
L_tu_j(t)=\mu_j(t)u_j(t),\qquad \mu_j(t)>0,
\]
and let $v_j(t)$ be the corresponding right singular vectors,
\[
J_tv_j(t)=\sqrt{\mu_j(t)}u_j(t),
\qquad
J_t^\ast u_j(t)=\sqrt{\mu_j(t)}v_j(t).
\]
Define residual target coefficients
\[
a_j(t)=\ip{e_t}{u_j(t)}.
\]

\begin{theorem}[Exact modal gradient energy]\label{thm:grad-energy}
At every checkpoint at which the singular system is defined,
\[
J_t^\ast e_t=\sum_j\sqrt{\mu_j(t)}a_j(t)v_j(t),
\]
and hence
\[
\boxed{
\norm{\nabla \R(\theta_t)}^2
=\norm{J_t^\ast e_t}^2
=\sum_j\mu_j(t)a_j(t)^2.}
\]
Moreover,
\[
\dot\R(\theta_t)=-\sum_j\mu_j(t)a_j(t)^2.
\]
\end{theorem}

\begin{proof}
Decompose $e_t=\Pi_te_t+(I-\Pi_t)e_t$, where $\Pi_t$ projects onto $\overline{\Ran J_t}$. Since $J_t^\ast(I-\Pi_t)e_t=0$,
\[
J_t^\ast e_t=\sum_j a_jJ_t^\ast u_j
=\sum_j\sqrt{\mu_j}a_jv_j.
\]
Orthonormality of the $v_j$ gives the norm identity. Finally,
\[
\dot\R=\ip{e_t}{\dot e_t}
=-\ip{e_t}{L_te_t}
=-\norm{J_t^\ast e_t}^2.
\]
\end{proof}

The product $\mu_j a_j^2$ is therefore not an arbitrary saliency heuristic. It is the exact instantaneous loss-dissipation budget of mode $j$.

\begin{proposition}[No eigenvalue-only safe deletion rule]\label{prop:no-eigen-only}
Fix any $0<\mu_2<\mu_1$. There exist two square-loss problems with the same prediction spectrum $\{\mu_1,\mu_2\}$ for which deleting the $\mu_2$ mode has, respectively, zero and strictly positive irreducible risk cost. Hence no rule depending only on eigenvalues can be uniformly safe.
\end{proposition}

\begin{proof}
Take $\Hh=\mathbb R^2$ and $L=\mathrm{diag}(\mu_1,\mu_2)$. In the first problem choose residual $e=(1,0)$; deleting the second mode changes nothing. In the second choose $e=(0,1)$; deleting the second mode prevents any reduction of the initial risk $1/2$. The spectrum is identical in both problems.
\end{proof}

\section{Fixed geometry: exact finite-horizon mode value}
We first solve the reduction problem exactly when $J_t\equiv J$ and $L_t\equiv L$ over a horizon $H$. This is the correct local benchmark for every nonlinear controller.

Let $S$ denote the kept singular modes and $D=S^c$ the dropped modes. Define the parameter projector
\[
Q_S=\sum_{j\in S}v_j\otimes v_j.
\]
The projected flow is
\[
\dot\theta=Q_SJ^\ast e,
\qquad
\dot e=-JQ_SJ^\ast e
=-\sum_{j\in S}\mu_j a_j u_j.
\]

\begin{theorem}[Exact finite-horizon deletion cost]\label{thm:fixed-value}
Suppose $L$ is fixed on $[0,H]$ and write
\[
e_0=e^\perp+\sum_j a_ju_j.
\]
The full and reduced terminal risks are
\begin{align*}
\R_{\rm full}(H)
&=\frac12\norm{e^\perp}^2
+\frac12\sum_j a_j^2e^{-2\mu_jH},\\
\R_S(H)
&=\frac12\norm{e^\perp}^2
+\frac12\sum_{j\in S}a_j^2e^{-2\mu_jH}
+\frac12\sum_{j\in D}a_j^2.
\end{align*}
Therefore the exact cost of deleting $D$ is
\[
\boxed{
\Delta_D(H)=\R_S(H)-\R_{\rm full}(H)
=\frac12\sum_{j\in D}a_j^2\bigl(1-e^{-2\mu_jH}\bigr).}
\]
For a single mode,
\[
V_j(H):=\frac12a_j^2(1-e^{-2\mu_jH})
\le \min\left\{H\mu_ja_j^2,\frac12a_j^2\right\},
\]
and as $\mu_jH\downarrow0$,
\[
V_j(H)=H\mu_ja_j^2+O(H^2\mu_j^2a_j^2).
\]
\end{theorem}

\begin{proof}
Each full mode solves $\dot a_j=-\mu_ja_j$, hence $a_j(H)=a_je^{-\mu_jH}$. Under the projected flow this remains true for $j\in S$, while $a_j$ is constant for $j\in D$. Orthogonality gives the risk formulas and their difference. The upper bound uses $1-e^{-x}\le\min\{x,1\}$ and the expansion follows from Taylor's theorem.
\end{proof}

This theorem corrects a common intuition. Small eigenvalues do not literally slow every other continuous-time mode. They are themselves slow, and they may waste compute; they affect global optimization rates only when the algorithm or guarantee depends on the least active curvature.

\subsection{An exact value-per-compute oracle}
Assign mode $j$ a risk-equivalent compute price $c_j\ge0$ per unit time. Over a horizon $H$, consider
\[
\mathcal J_H(S)=\R_S(H)+H\sum_{j\in S}c_j.
\]

\begin{corollary}[Oracle active set]\label{cor:value-compute}
Under the assumptions of Theorem~\ref{thm:fixed-value}, an optimal active set is obtained independently mode by mode:
\[
\boxed{j\in S_H^\star\quad\Longleftrightarrow\quad V_j(H)>Hc_j,}
\]
with either decision admissible at equality. For $\mu_jH\ll1$, the criterion reduces to
\[
\mu_ja_j^2\gtrsim c_j.
\]
\end{corollary}

\begin{proof}
The objective is a sum of modewise terms plus constants. Keeping $j$ decreases terminal risk by $V_j(H)$ and costs $Hc_j$.
\end{proof}

\section{Optimization consequence of truncation}
Consider the fixed linearized quadratic problem restricted to a kept set $S$ with
\[
0<\mu_{\min,S}\le\mu_j\le\mu_{\max,S}.
\]

\begin{proposition}[Conditioning of the active problem]\label{prop:conditioning}
Gradient descent on the kept quadratic subproblem with scalar step
\[
\eta_S^\star=\frac{2}{\mu_{\max,S}+\mu_{\min,S}}
\]
has worst-case contraction factor
\[
\rho_S=\frac{\kappa_S-1}{\kappa_S+1},
\qquad
\kappa_S=\frac{\mu_{\max,S}}{\mu_{\min,S}}.
\]
Removing a tail below a threshold $\tau$ guarantees $\mu_{\min,S}\ge\tau$ and therefore improves this condition-number bound whenever the removed tail determined the previous minimum eigenvalue.
\end{proposition}

\begin{proof}
On each kept mode the error multiplier is $1-\eta\mu_j$. Minimizing the maximum absolute multiplier over the interval $[\mu_{\min,S},\mu_{\max,S}]$ equalizes the two endpoint magnitudes, yielding the stated step and factor.
\end{proof}

\begin{remark}
This is a statement about the conditioned active problem, not a universal promise that pruning accelerates dense gradient descent. If the same stable step is used in full and reduced models, deleting a weak mode does not make an unrelated large-$\mu$ mode decay faster. Practical acceleration requires either fewer structural FLOPs, a step/preconditioner that exploits the improved active spectrum, or both.
\end{remark}

\section{Finite-sample statistical value of deleting a mode}
The second possible gain is statistical rather than purely computational. In the Gaussian sequence model, observe
\[
z_j=a_j+\frac{\sigma}{\sqrt n}\xi_j,
\qquad \xi_j\sim N(0,1),
\]
and use ridge filter $q_\lambda(\mu)=\mu/(\mu+\lambda)$. The risk contribution of an active mode is
\[
R_j^{\rm keep}
=\left(\frac{\lambda}{\mu_j+\lambda}\right)^2a_j^2
+\frac{\sigma^2}{n}\left(\frac{\mu_j}{\mu_j+\lambda}\right)^2.
\]
Dropping the mode sets its estimate to zero and gives $R_j^{\rm drop}=a_j^2$.

\begin{theorem}[Exact finite-sample drop criterion]\label{thm:sequence-drop}
For $\mu_j>0$,
\[
\boxed{
R_j^{\rm drop}-R_j^{\rm keep}
=
\frac{\mu_j}{(\mu_j+\lambda)^2}
\left[a_j^2(\mu_j+2\lambda)-\frac{\sigma^2}{n}\mu_j\right].}
\]
Hence dropping the mode improves conditional sequence risk exactly when
\[
\boxed{
a_j^2(\mu_j+2\lambda)<\frac{\sigma^2}{n}\mu_j.}
\]
At $\lambda=0$, this becomes $a_j^2<\sigma^2/n$.
\end{theorem}

\begin{proof}
Subtract the two displayed risks and use
\[
1-\frac{\lambda^2}{(\mu+\lambda)^2}
=\frac{\mu(\mu+2\lambda)}{(\mu+\lambda)^2}.
\]
\end{proof}

Thus nuisance directions can be harmful even if they are learnable: they may contribute less target signal than sampling variance. This is distinct from the finite-horizon optimization value in Theorem~\ref{thm:fixed-value}; both should enter a compute-aware decision.

\section{Compute-to-accuracy after spectral contraction}
The conditioning proposition can be sharpened into a direct compute statement. The important qualification is that deleting modes creates an approximation floor. Acceleration is meaningful only for a target accuracy that remains above that floor.

Keep a fixed set $S$ and drop $D=S^c$. Write
\[
F_\perp:=\frac12\norm{e^\perp}^2,
\qquad
F_D:=\frac12\sum_{j\in D}a_j^2,
\qquad
E_S(0):=\frac12\sum_{j\in S}a_j^2.
\]
For scalar-step gradient descent on the retained fixed quadratic problem, use the optimal step
\[
\eta_S^\star=\frac{2}{\mu_{\max,S}+\mu_{\min,S}},
\qquad
\rho_S=\frac{\kappa_S-1}{\kappa_S+1}.
\]

\begin{theorem}[Iteration complexity above the deletion floor]\label{thm:compute-accuracy}
For every integer $k\ge0$,
\[
\boxed{
\R_S(k)-F_\perp-F_D
\le
\rho_S^{2k}E_S(0).}
\]
Hence for any absolute target risk $R_{\rm tar}>F_\perp+F_D$, it is sufficient that
\[
\boxed{
k\ge
K_S(R_{\rm tar})
:=
\left\lceil
\frac{\log\!\left(E_S(0)/(R_{\rm tar}-F_\perp-F_D)\right)_+}
{-2\log\rho_S}
\right\rceil,}
\]
with the convention $K_S=0$ when the numerator is nonpositive. If one retained iteration costs $C_S$, the corresponding compute bound is
\[
\boxed{\mathcal C_S(R_{\rm tar})\le C_SK_S(R_{\rm tar}).}
\]
For the dense active set replace $S$ by all positive modes and set $F_D=0$.
\end{theorem}

\begin{proof}
Every retained modal coefficient is multiplied by $1-\eta_S^\star\mu_j$, whose absolute value is at most $\rho_S$. Squaring and summing the orthogonal modal risk contributions gives the first inequality. Solving
$\rho_S^{2k}E_S(0)\le R_{\rm tar}-F_\perp-F_D$
for $k$ gives the second statement. Multiplication by the per-iteration cost gives the compute bound.
\end{proof}

\begin{corollary}[When spectral contraction improves the theoretical compute bound]\label{cor:compute-improve}
Let $K_{\rm full}(R_{\rm tar})$ and $C_{\rm full}$ denote the corresponding dense quantities. Whenever
\[
F_D<R_{\rm tar}-F_\perp
\]
and
\[
\boxed{C_SK_S(R_{\rm tar})<C_{\rm full}K_{\rm full}(R_{\rm tar}),}
\]
the reduced problem has a strictly smaller certified compute-to-target bound. This improvement can arise from either or both of
\[
C_S<C_{\rm full}
\qquad\text{and}\qquad
\kappa_S<\kappa_{\rm full}.
\]
\end{corollary}

\begin{proof}
The floor condition makes the target attainable by the reduced model. The second inequality directly compares the sufficient compute bounds from Theorem~\ref{thm:compute-accuracy}.
\end{proof}

\begin{remark}[The precise version of ``small eigenvalues slow training'']
In continuous-time gradient flow a weak mode does not slow an unrelated strong mode. In scalar-step discrete gradient descent, however, a very small retained eigenvalue can worsen the global condition number and therefore the worst-case iteration bound. If that mode also has negligible finite-horizon target value, deleting it can improve the condition number without materially changing the attainable risk. Structural deletion can then add a second gain by lowering $C_S$. The theorem makes both effects explicit and keeps the approximation floor visible.
\end{remark}

\section{When temporary overparameterization can pay for itself}
The previous compute-to-accuracy result asks whether a contraction is beneficial relative to continuing the same wide state. A harder question is the one relevant for the claim that a network should \emph{learn wide and then become small}: why not train the final compact architecture from the beginning? The answer cannot be ``because contraction is safe.'' Safety only says that little is lost at the contraction event. Temporary width is useful at equal compute only if the wide phase creates a sufficiently better compact state to repay the compute spent while wide.

Fix a wide architecture $W$ and a realizable compact architecture $r$, with per-update costs $C_W\ge C_r>0$. Starting from a common target-free initialization rule, let
\[
R_{W\to r}(t,H)
\]
denote the risk obtained by training the wide model for $t$ updates, structurally contracting it to architecture $r$, and then training the compact model for another $H$ updates. Let $R_r(k)$ denote the risk of the compact-from-start control after $k$ updates. Define the equal-compute compact budget
\[
N_r(t,H)
:=
\left\lceil\frac{C_Wt+C_rH}{C_r}\right\rceil,
\qquad
K(t,H):=t+H.
\]
The ceiling weakly favors the compact control and can be dropped when the compute ratio is exact.

\begin{definition}[Discovery advantage and compute opportunity gain]
Define
\[
D_r(t,H)
:=R_r(K(t,H))-R_{W\to r}(t,H)
\]
and
\[
O_r(t,H)
:=R_r(K(t,H))-R_r(N_r(t,H)).
\]
$D_r$ measures how much better the state produced by temporary width is than the compact-from-start state at the \emph{same update count}. It includes every useful effect of temporary overparameterization---representation discovery, optimization-path selection, conditioning, and initialization amplification---after both paths have been expressed in the same final architecture. $O_r$ measures how much the compact control gains from the additional updates it can buy with the compute that the wide phase consumed.
\end{definition}

\begin{theorem}[Exact break-even identity]\label{thm:temporary-width-break-even}
For every $t,H\ge0$,
\[
\boxed{
R_{W\to r}(t,H)-R_r(N_r(t,H))
=
O_r(t,H)-D_r(t,H).}
\]
Consequently, temporary overparameterization wins the equal-compute comparison if and only if
\[
\boxed{D_r(t,H)>O_r(t,H).}
\]
If $R_r(k)$ is nonincreasing in $k$ and $C_W\ge C_r$, then $O_r(t,H)\ge0$. In that case $D_r(t,H)>0$ is a necessary condition for any equal-compute advantage.
\end{theorem}

\begin{proof}
Add and subtract $R_r(K(t,H))$:
\begin{align*}
R_{W\to r}(t,H)-R_r(N_r(t,H))
&=
\bigl[R_{W\to r}(t,H)-R_r(K(t,H))\bigr]\\
&\quad+
\bigl[R_r(K(t,H))-R_r(N_r(t,H))\bigr]\\
&=-D_r(t,H)+O_r(t,H).
\end{align*}
The first equivalence follows immediately. Moreover $C_W\ge C_r$ implies $N_r(t,H)\ge K(t,H)$; monotonicity of the compact risk then gives $R_r(K)\ge R_r(N_r)$ and hence $O_r\ge0$.
\end{proof}

\begin{corollary}[Safe contraction is not enough]\label{cor:safety-not-enough}
Suppose the structural contraction at time $t$ has arbitrarily small finite-horizon damage relative to continuing the wide model. This fact alone places no sign restriction on $D_r(t,H)$ and therefore cannot imply an equal-compute advantage over compact-from-start. In particular, if $D_r(t,H)\le0$ and compact risk is nonincreasing, temporary overparameterization cannot win at equal compute.
\end{corollary}

\begin{proof}
The contraction-damage comparison uses a continuation of the checkpointed wide state as its reference, whereas $D_r$ compares against an independently trained compact path. These are different counterfactuals. The final statement follows from Theorem~\ref{thm:temporary-width-break-even} because $O_r\ge0$.
\end{proof}

\begin{remark}[The empirical decomposition that should be reported]
An equal-compute experiment should report all three quantities
\[
R_{W\to r}(t,H)-R_r(K),
\qquad
O_r(t,H),
\qquad
R_{W\to r}(t,H)-R_r(N_r),
\]
rather than only the last one. The first term reveals whether temporary width actually discovered a better compact state; the second is the price of having trained wide; the third is their net effect. This decomposition distinguishes a failed discovery mechanism from a useful discovery mechanism that is simply too expensive.
\end{remark}

\begin{remark}[Research consequence]
The break-even theorem changes the role of the controller. A safe-pruning certificate determines when one \emph{may} contract. A useful learn-wide-then-shrink algorithm needs an additional economic gate: estimated discovery advantage must be large enough to have justified the historical wide-phase compute, or, prospectively, the expected future discovery value of remaining wide must exceed its incremental compute price. This is stronger than spectral truncation and stronger than phase separation alone.
\end{remark}

\section{A local compute-efficiency threshold}
The break-even identity is retrospective: it asks whether the wide phase has already created enough discovery advantage to repay its historical compute. For an online controller we also need a prospective question: starting from the same current predictor, is the next infinitesimal unit of compute better spent in the wide parameter space or in a reduced subspace? In the ideal spectral model this comparison is exact and produces a dimensionless threshold.

Let $g=\nabla_\theta\R(\theta)$ at the current checkpoint and let $Q$ be an orthogonal projector onto the parameter directions retained by a candidate contraction. Consider two continuous-time updates starting from the same parameter point,
\[
\dot\theta_W=-g(\theta_W),
\qquad
\dot\theta_Q=-Q g(\theta_Q).
\]
Suppose one unit of gradient-flow time costs $C_W$ for the full update and $C_Q$ for the projected update, with $0<C_Q\le C_W$. Parameterize both trajectories by cumulative compute $b$, so that an increment $db$ corresponds locally to $db/C_W$ units of wide flow or $db/C_Q$ units of projected flow.

\begin{theorem}[Exact infinitesimal value per unit compute]\label{thm:local-compute-efficiency}
At the common checkpoint,
\[
\left.\frac{d}{db}\R(\theta_W(b))\right|_{b=0}
=-\frac{\norm{g}^2}{C_W},
\qquad
\left.\frac{d}{db}\R(\theta_Q(b))\right|_{b=0}
=-\frac{\norm{Qg}^2}{C_Q}.
\]
Therefore the reduced update has strictly larger instantaneous risk decrease per unit compute if and only if
\[
\boxed{
\frac{\norm{Qg}^2}{C_Q}
>
\frac{\norm{g}^2}{C_W}.}
\]
Writing
\[
A_S:=\norm{Qg}^2,
\qquad
A_D:=\norm{(I-Q)g}^2,
\qquad
A=A_S+A_D,
\]
this condition is equivalently
\[
\boxed{
A_D
<
\left(\frac{C_W}{C_Q}-1\right)A_S,}
\]
or, whenever $A>0$,
\[
\boxed{
\frac{A_D}{A}
<
1-\frac{C_Q}{C_W}.}
\]
Thus the fraction of instantaneous gradient-descent power discarded by contraction must be smaller than the fraction of per-update compute saved.
\end{theorem}

\begin{proof}
Along ordinary gradient flow,
\[
\frac{d}{ds}\R(\theta_W(s))
=
\ip{g}{-g}
=-\norm{g}^2.
\]
Since $s=b/C_W$, the chain rule gives the first derivative. Along the projected flow,
\[
\frac{d}{ds}\R(\theta_Q(s))
=
\ip{g}{-Qg}
=-\norm{Qg}^2,
\]
because $Q$ is an orthogonal projector; reparameterizing by $s=b/C_Q$ gives the second derivative. Since $Q$ and $I-Q$ have orthogonal ranges,
$\norm{g}^2=A_S+A_D$. The reduced update is locally more compute-efficient exactly when
$A_S/C_Q>(A_S+A_D)/C_W$. Rearranging gives the second boxed inequality. Dividing by $A$ gives the last one.
\end{proof}

\begin{corollary}[Target-weighted spectral form]\label{cor:local-compute-spectral}
Let $J=U\Sigma V^\ast$, let $\mu_j$ be the positive eigenvalues of $JJ^\ast$, and let $a_j=\ip{e}{u_j}$. If $Q$ retains the right singular directions indexed by $S$ and drops $D=S^c$, then
\[
A_S=\sum_{j\in S}\mu_j a_j^2,
\qquad
A_D=\sum_{j\in D}\mu_j a_j^2.
\]
Consequently the candidate is instantaneously more efficient per unit compute if and only if
\[
\boxed{
\frac{\sum_{j\in D}\mu_j a_j^2}
{\sum_j\mu_j a_j^2}
<
1-\frac{C_Q}{C_W}.}
\]
\end{corollary}

\begin{proof}
The modal gradient identity gives
$g=J^\ast e=\sum_j\sqrt{\mu_j}a_jv_j$.
Orthogonality of the $v_j$ yields the two energy sums. Apply Theorem~\ref{thm:local-compute-efficiency}.
\end{proof}

\begin{remark}[Why this is not yet a pruning theorem]
The inequality above answers only the next-infinitesimal-compute question. It does not control the structural jump produced by a physical contraction, the approximation floor created by permanently deleting target mass, or the possibility that currently weak directions later become target-relevant because the Jacobian rotates. Those are controlled by the structural-transfer, finite-horizon, and future-discovery certificates developed elsewhere in this note. The useful interpretation is therefore not ``prune whenever the inequality holds,'' but rather: among candidates that are already statistically and geometrically safe, reject any contraction that cannot even beat the wide model in instantaneous risk decrease per unit compute.
\end{remark}

\begin{remark}[A dimensionless onset statistic]
Define the omitted descent-power fraction
\[
\omega_D(t):=
\frac{\norm{(I-Q_t)\nabla\R(\theta_t)}^2}
{\norm{\nabla\R(\theta_t)}^2}
\]
and the structural compute-saving fraction
\[
s_C(t):=1-\frac{C_{Q_t}}{C_W}.
\]
The local compute gate is simply $\omega_D(t)<s_C(t)$. Both sides are dimensionless. The left side is target-weighted prediction geometry; the right side is a property of the actual structural implementation. This is the cleanest local bridge between the spectral theory and physical compute.
\end{remark}

\begin{corollary}[First-order break-even over a small compute budget]\label{cor:local-break-even}
Assume the two trajectories above are differentiable in compute at $b=0$. For a small equal-compute budget $B>0$,
\[
\R(\theta_Q(B))-\R(\theta_W(B))
=
B\left(
\frac{\norm{g}^2}{C_W}
-
\frac{\norm{Qg}^2}{C_Q}
\right)+o(B).
\]
Hence the sign in Theorem~\ref{thm:local-compute-efficiency} is the exact first-order sign of the equal-compute risk difference.
\end{corollary}

\begin{proof}
Apply the first-order Taylor expansion in the common compute parameter $b$ to both risk trajectories and subtract.
\end{proof}

\section{Exact finite-budget equal-compute contraction in fixed geometry}
The infinitesimal compute threshold has an exact finite-budget analogue whenever the prediction geometry is frozen. This gives the cleanest possible prospective decision rule: compare the learning progress forfeited in deleted modes with the additional progress purchased in retained modes by making each update cheaper.

Fix a constant prediction operator
\[
L=\sum_j\mu_j u_ju_j^\ast,
\qquad
 e_0=\sum_j a_ju_j+e^\perp,
\]
and a retained index set $S$ with deleted complement $D$. Let one unit of full gradient-flow time cost $C_W>0$ and one unit of reduced flow time cost $C_S\in(0,C_W]$. For an equal cumulative compute budget $B\ge0$, the full path receives flow time
\[
h_W=\frac{B}{C_W},
\]
whereas the reduced path receives
\[
h_S=\frac{B}{C_S}\ge h_W.
\]
The reduced path evolves the modes in $S$ and freezes the modes in $D$.

\begin{theorem}[Exact finite-budget compute decomposition]\label{thm:fixed-equal-compute}
At equal compute $B$,
\[
\R_W(B)
=
\frac12\sum_j a_j^2e^{-2\mu_jh_W}
+\frac12\norm{e^\perp}^2,
\]
while
\[
\R_S(B)
=
\frac12\sum_{j\in S}a_j^2e^{-2\mu_jh_S}
+\frac12\sum_{j\in D}a_j^2
+\frac12\norm{e^\perp}^2.
\]
Define the deleted-mode opportunity loss
\[
L_D(B)
:=
\frac12\sum_{j\in D}a_j^2
\left(1-e^{-2\mu_jB/C_W}\right)
\]
and the retained-mode acceleration gain
\[
G_S(B)
:=
\frac12\sum_{j\in S}a_j^2
\left(
 e^{-2\mu_jB/C_W}
-e^{-2\mu_jB/C_S}
\right).
\]
Then
\[
\boxed{
\R_S(B)-\R_W(B)=L_D(B)-G_S(B).}
\]
Consequently the contraction strictly improves risk at the same compute budget if and only if
\[
\boxed{G_S(B)>L_D(B).}
\]
Both terms are nonnegative. The left side is exactly the extra progress on retained modes bought by the cheaper reduced dynamics; the right side is exactly the progress the full model would have made on the modes that were deleted.
\end{theorem}

\begin{proof}
Because $L$ is fixed, every full-flow modal coefficient evolves as
$a_j(s)=a_je^{-\mu_js}$. Evaluating at $s=h_W$ gives the first risk expression. Under the reduced dynamics, retained coefficients evolve for the longer equal-compute time $h_S$, while deleted coefficients remain equal to their initial values. This gives the second expression. Subtracting the two risks and separating $S$ from $D$ yields
\begin{align*}
\R_S(B)-\R_W(B)
&=
\frac12\sum_{j\in S}a_j^2
\left(e^{-2\mu_jh_S}-e^{-2\mu_jh_W}\right)\\
&\quad+
\frac12\sum_{j\in D}a_j^2
\left(1-e^{-2\mu_jh_W}\right)\\
&=-G_S(B)+L_D(B).
\end{align*}
Since $h_S\ge h_W$ and $\mu_j\ge0$, both definitions are nonnegative. The final equivalence is immediate.
\end{proof}

\begin{corollary}[The local compute threshold is the small-budget limit]\label{cor:fixed-to-local}
As $B\downarrow0$,
\[
L_D(B)
=
\frac{B}{C_W}\sum_{j\in D}\mu_ja_j^2+O(B^2),
\]
and
\[
G_S(B)
=
B\left(\frac1{C_S}-\frac1{C_W}\right)
\sum_{j\in S}\mu_ja_j^2+O(B^2).
\]
Therefore $G_S(B)>L_D(B)$ for all sufficiently small positive $B$ exactly when
\[
\sum_{j\in D}\mu_ja_j^2
<
\left(\frac{C_W}{C_S}-1\right)
\sum_{j\in S}\mu_ja_j^2,
\]
which is Corollary~\ref{cor:local-compute-spectral}.
\end{corollary}

\begin{proof}
Use $1-e^{-x}=x+O(x^2)$ and
$e^{-x}-e^{-y}=(y-x)+O(x^2+y^2)$ for $x,y\to0$, then simplify.
\end{proof}

\begin{remark}[Why eigenvalue thresholding is insufficient even with a compute budget]
The decision depends on $a_j^2$ and on the time horizon $B$, not on $\mu_j$ alone. A small-$\mu_j$ direction with large target mass can have substantial finite-budget opportunity loss, while a larger eigenvalue carrying essentially no residual target mass can be cheap to delete. The compute ratio also matters: an aggressive structural contraction can justify a larger deleted-mode loss because it buys more evolution time for every retained mode.
\end{remark}

\begin{remark}[The role of moving geometry]
The theorem is exact but deliberately frozen-geometry. In a nonlinear network, a deleted sector can acquire future target value because $J_t$ and its eigenspaces move. The dynamic safety results therefore add an explicit geometry-motion allowance to $L_D(B)$ before using this comparison as a contraction certificate. Conceptually the nonlinear controller should preserve the same accounting identity: \emph{future value forgone by deletion} versus \emph{future progress and compute saved by contraction}.
\end{remark}

\section{Moving geometry: why pruning can be premature}
In a feature-learning network, the basis itself moves. For simple eigenvalues,
\[
\dot\mu_j=\ip{u_j}{\dot L_tu_j},
\]
and under the gauge $\ip{u_j}{\dot u_j}=0$,
\[
\dot u_j
=\sum_{m\ne j}
\frac{\ip{u_m}{\dot L_tu_j}}{\mu_j-\mu_m}u_m.
\]
The residual coefficients satisfy
\[
\boxed{
\dot a_j
=-\mu_ja_j
+\sum_{m\ne j}
\frac{\ip{u_m}{\dot L_tu_j}}{\mu_j-\mu_m}a_m.}
\]
The first term is coefficient fitting in the current basis; the second is target transport caused by feature learning. A currently weak mode can therefore become important because its eigenvalue grows, its eigenspace rotates toward residual target mass, or both.

This fact motivates a \emph{persistence} requirement: deletion should be based on a finite-horizon upper bound, not a single checkpoint.

\section{Curvature controls how fast prediction geometry can move}
The network does not move its prediction operator arbitrarily fast. The exact gradient-dissipation identity already limits the distance that the geometry can travel when curvature is controlled.

For the second Fr\'echet derivative define
\[
\norm{D^2F(\theta)}
:=\sup_{\norm v=\norm w=1}\norm{D^2F(\theta)[v,w]}_{\Hh}.
\]

\begin{theorem}[Curvature-to-spectral-motion bound]\label{thm:curvature-motion}
Along square-loss gradient flow, assume on $[t,t+H]$ that
\[
\norm{J_s}_{\op}\le B_J,
\qquad
\norm{D^2F(\theta_s)}\le B_H.
\]
Then
\[
\boxed{
\norm{\dot L_s}_{\op}
\le 2B_HB_J\norm{J_s^\ast e_s}.}
\]
Consequently the operator path length satisfies
\[
\boxed{
\int_t^{t+H}\norm{\dot L_s}_{\op}\,ds
\le
2B_HB_J\sqrt{H\,[\R(\theta_t)-\R(\theta_{t+H})]}
\le
2B_HB_J\sqrt{H\R(\theta_t)}.}
\]
If $P_I(s)$ is the Riesz projector of a spectral cluster separated from the rest of $L_s$ by a gap at least $\Delta>0$, then
\[
\boxed{
\int_t^{t+H}\norm{\dot P_I(s)}_{\op}\,ds
\le
\frac{4B_HB_J}{\Delta}
\sqrt{H\,[\R(\theta_t)-\R(\theta_{t+H})]}.}
\]
\end{theorem}

\begin{proof}
Differentiating $J_s=DF(\theta_s)$ along the trajectory gives
\[
\dot J_s[v]=D^2F(\theta_s)[\dot\theta_s,v],
\]
so
\[
\norm{\dot J_s}_{\op}\le B_H\norm{\dot\theta_s}
=B_H\norm{J_s^\ast e_s}.
\]
Since $L_s=J_sJ_s^\ast$,
\[
\dot L_s=\dot J_sJ_s^\ast+J_s\dot J_s^\ast,
\]
and hence
\[
\norm{\dot L_s}_{\op}
\le2\norm{J_s}_{\op}\norm{\dot J_s}_{\op}
\le2B_HB_J\norm{J_s^\ast e_s}.
\]
By Theorem~\ref{thm:grad-energy},
\[
\int_t^{t+H}\norm{J_s^\ast e_s}^2ds
=\R(\theta_t)-\R(\theta_{t+H}).
\]
Cauchy--Schwarz proves the operator path-length bound. Finally, the standard differentiable Riesz-projector bound
$\norm{\dot P_I}_{\op}\le2\norm{\dot L}_{\op}/\Delta$
gives the cluster path-length claim.
\end{proof}

\begin{remark}
This result supplies the missing dynamical input to a feature-value stopping argument. A large unresolved target component is not enough to justify continued full-network training: the network must also have enough remaining curvature-driven subspace motion to reach it. Conversely, early in training a large geometry path budget is a mathematical reason not to prune merely because current eigenvalues are small.
\end{remark}

\section{Finite-horizon stability of structural contraction}
We now compare nonlinear full training with a physically reduced flow. Define
\[
g(\theta)=J_\theta^\ast e_\theta=-\nabla\R(\theta).
\]
Let $S$ be a fixed orthogonal projector representing the retained structural coordinates after contraction. We allow an instantaneous approximation error at the contraction checkpoint, because a reduced architecture need not be an exact embedding of the dense state. Starting at time $t$, write
\[
\dot\theta_s=g(\theta_s),
\qquad
\dot{\widetilde\theta}_s=Sg(\widetilde\theta_s),
\qquad
\norm{\theta_t-\widetilde\theta_t}=d_0.
\]
The reduced parameters may be embedded in the original parameter space by setting deleted structural coordinates inactive.

\begin{theorem}[Safe finite-horizon structural-flow bound]\label{thm:gronwall}
Assume on a tube containing both trajectories for $s\in[t,t+H]$ that:
\begin{enumerate}[label=(\roman*)]
\item $g$ is $L_g$-Lipschitz;
\item $\norm{J_\theta}_{\op}\le B_J$;
\item both residual norms are at most $B_e$;
\item the omitted dense-path gradient obeys
$\norm{(I-S)g(\theta_s)}\le\varepsilon(s)$.
\end{enumerate}
Then
\[
\boxed{
\norm{\theta_{t+H}-\widetilde\theta_{t+H}}
\le
e^{L_gH}d_0
+\int_t^{t+H}e^{L_g(t+H-s)}\varepsilon(s)\,ds,}
\]
\[
\boxed{
\norm{F(\theta_{t+H})-F(\widetilde\theta_{t+H})}
\le
B_J\left[
e^{L_gH}d_0
+\int_t^{t+H}e^{L_g(t+H-s)}\varepsilon(s)\,ds
\right],}
\]
and
\[
\boxed{
|\R(\theta_{t+H})-\R(\widetilde\theta_{t+H})|
\le
B_eB_J\left[
e^{L_gH}d_0
+\int_t^{t+H}e^{L_g(t+H-s)}\varepsilon(s)\,ds
\right].}
\]
If $d_0=0$ and $\varepsilon(s)\le\bar\varepsilon$, the dynamic term is
$\bar\varepsilon(e^{L_gH}-1)/L_g$, with continuous limit $H\bar\varepsilon$ when $L_g=0$.
\end{theorem}

\begin{proof}
Set $z_s=\theta_s-\widetilde\theta_s$. Since $S$ is a contraction,
\[
\dot z_s
=(I-S)g(\theta_s)+S[g(\theta_s)-g(\widetilde\theta_s)],
\]
so
\[
\frac{d}{ds}\norm{z_s}
\le\varepsilon(s)+L_g\norm{z_s}
\]
almost everywhere. Gr\"onwall's inequality gives the parameter bound including the initial discrepancy $d_0$. The Jacobian bound makes $F$ $B_J$-Lipschitz on the tube. Finally,
\[
|\R(\theta)-\R(\tilde\theta)|
\le\frac12(\norm{e_\theta}+\norm{e_{\tilde\theta}})
\norm{e_\theta-e_{\tilde\theta}}
\le B_e\norm{F(\theta)-F(\tilde\theta)}.
\]
\end{proof}

This formulation is useful operationally because the omitted-gradient certificate may be estimated or bounded along the dense reference trajectory before an irreversible contraction is executed.

\section{A spectral persistence certificate}
Let $P_D(t)$ be the Riesz projector onto an isolated positive spectral cluster $D$ of $L_t$, separated from the rest of the spectrum by a gap $\Delta_t\ge\Delta>0$. Define the cluster gradient energy
\[
E_D(t)=\norm{J_t^\ast P_D(t)e_t}^2
=\ip{e_t}{L_tP_D(t)e_t}.
\]
For a simple spectrum this is $\sum_{j\in D}\mu_j a_j^2$.

\begin{theorem}[Upper growth bound for discarded-cluster gradient energy]\label{thm:persistence}
Suppose $L_t$ is differentiable, $P_D(t)$ remains isolated by gap $\Delta$, and square-loss gradient flow holds. Then
\[
\dot E_D(t)
\le
\norm{e_t}^2\norm{\dot L_t}_{\op}
\left(1+\frac{2\norm{L_t}_{\op}}{\Delta}\right).
\]
If on $[t,t+H]$
\[
\norm{e_s}\le B_e,
\qquad
\norm{L_s}_{\op}\le B_L,
\]
then for $0\le h\le H$,
\[
\boxed{
E_D(t+h)
\le E_D(t)
+B_e^2\left(1+\frac{2B_L}{\Delta}\right)
\int_t^{t+h}\norm{\dot L_s}_{\op}\,ds.}
\]
\end{theorem}

\begin{proof}
Because $P_D$ is a spectral projector, it commutes with $L$. Put $A=LP_D$. Then
\[
E_D=\ip e{Ae},
\quad
\dot E_D=2\ip{\dot e}{Ae}+\ip e{\dot A e}.
\]
Using $\dot e=-Le$ and $LP_D=P_DL$,
\[
2\ip{\dot e}{Ae}
=-2\ip{Le}{LP_De}
=-2\norm{LP_De}^2\le0.
\]
Also $\dot A=\dot L P_D+L\dot P_D$, so
\[
\dot E_D\le\norm e^2
\left(\norm{\dot L}_{\op}+\norm L_{\op}\norm{\dot P_D}_{\op}\right).
\]
The Riesz-projector/Davis--Kahan derivative bound gives
$\norm{\dot P_D}_{\op}\le2\norm{\dot L}_{\op}/\Delta$. Integrating proves the claim.
\end{proof}

\begin{corollary}[Curvature-based finite-horizon omitted-energy envelope]\label{cor:energy-envelope}
Under Theorems~\ref{thm:curvature-motion} and~\ref{thm:persistence}, define
\[
C_D=B_e^2\left(1+\frac{2B_L}{\Delta}\right),
\qquad
V_L(t,H)=2B_HB_J\sqrt{H[\R(\theta_t)-\R(\theta_{t+H})]}.
\]
Then
\[
E_D(t+h)\le E_D(t)+C_DV_L(t,H),
\qquad 0\le h\le H,
\]
and
\[
\boxed{
\int_t^{t+H}E_D(s)\,ds
\le H\,[E_D(t)+C_DV_L(t,H)].}
\]
In particular,
\[
\int_t^{t+H}\sqrt{E_D(s)}\,ds
\le
H\sqrt{E_D(t)+C_DV_L(t,H)}.
\]
\end{corollary}

The gap is essential. Inside an unresolved random-matrix bulk, individual empirical eigenvectors are unstable; deletion decisions should be made at the cluster level until a stable subspace detaches.

\section{From ideal spectral reduction to actual parameter deletion}
A spectral projector is not automatically a hardware-friendly parameter mask. This is the structural-realizability gap.

Let $Q_K$ be an ideal kept parameter projector and let $S$ be a structural projector induced by a realizable architecture reduction (for example, retained neurons or channels).

\begin{theorem}[Structural transfer]\label{thm:structural-transfer}
At any checkpoint,
\[
\boxed{
\norm{(I-S)g}
\le
\norm{(I-Q_K)g}
+\norm{(I-S)Q_Kg}.}
\]
In particular, if the ideal dropped sector has energy
$\norm{(I-Q_K)g}^2\le E_D$ and the structural approximation satisfies
$\norm{(I-S)Q_Kg}\le\delta$, then
\[
\norm{(I-S)g}\le\sqrt{E_D}+\delta.
\]
\end{theorem}

\begin{proof}
Write
$(I-S)g=(I-S)(I-Q_K)g+(I-S)Q_Kg$ and use that an orthogonal projector has operator norm one.
\end{proof}

The pointwise theorem is enough for a checkpoint. For future safety, however, a structural mask chosen at time $t$ must continue to approximate the moving kept spectral subspace.

\begin{theorem}[Persistence of a fixed structural realization]\label{thm:structural-persistence}
Let $G_s=J_s^\ast J_s$ be the parameter Gram operator. Suppose its ideal kept cluster has spectral projector $Q_K(s)$ and remains separated by a gap at least $\Delta_G>0$ on $[t,t+H]$. Let $S$ be a fixed structural projector selected at time $t$, and define
\[
\delta_0=\norm{(I-S)Q_K(t)}_{\op}.
\]
Under the curvature assumptions of Theorem~\ref{thm:curvature-motion}, for $0\le h\le H$,
\[
\boxed{
\norm{(I-S)Q_K(t+h)}_{\op}
\le
\delta_0
+\frac{4B_HB_J}{\Delta_G}
\sqrt{h\,[\R(\theta_t)-\R(\theta_{t+h})]}.}
\]
Consequently, with $Q_D(s)=I-Q_K(s)$,
\[
\boxed{
\norm{(I-S)g(\theta_s)}
\le
\norm{Q_D(s)g(\theta_s)}
+\delta_S(s)\norm{g(\theta_s)},}
\]
where $\delta_S(s)$ is bounded by the preceding display.
\end{theorem}

\begin{proof}
As for $L_s$, differentiating $G_s=J_s^\ast J_s$ gives
\[
\norm{\dot G_s}_{\op}
\le2B_HB_J\norm{J_s^\ast e_s}.
\]
The Riesz-projector derivative bound in parameter space yields
\[
\norm{Q_K(t+h)-Q_K(t)}_{\op}
\le\frac{2}{\Delta_G}\int_t^{t+h}\norm{\dot G_s}_{\op}ds.
\]
Apply Cauchy--Schwarz and the exact loss-dissipation identity to obtain the first claim. Finally,
\[
(I-S)g=(I-S)Q_Dg+(I-S)Q_Kg
\]
and the projector contraction bound gives the second claim.
\end{proof}

\begin{corollary}[Modular safe-contraction certificate]\label{cor:master-certificate}
A sufficient finite-horizon certificate for a structural contraction is obtained by the following chain:
\[
\begin{gathered}
\text{curvature + current loss}
\Rightarrow
\text{spectral path budget}
\Rightarrow
\text{discarded-energy persistence}\\
\Downarrow\\
\text{structural omitted-gradient budget}
\Rightarrow
\text{terminal risk bound}
\end{gathered}
\]
Specifically, Theorem~\ref{thm:curvature-motion} bounds future spectral motion, Theorem~\ref{thm:persistence} bounds the ideal discarded energy, Theorem~\ref{thm:structural-persistence} controls a fixed structural realization, and Theorem~\ref{thm:gronwall} converts the resulting $\varepsilon(s)$ together with any initial contraction distortion $d_0$ into an explicit terminal prediction/risk error.
\end{corollary}

This forces an honest separation between two scientific claims: prediction geometry may say that a sector is expendable, while a particular neuron/channel realization may still approximate that sector poorly. The structural gap is not a nuisance to hide; it is part of the theorem.

\section{Exact current sufficiency for the final hidden layer}
The final hidden layer followed by a scalar linear readout gives the cleanest structural case. Let
\[
h_t:\mathcal X\to\mathbb R^m,
\qquad
(T_t\beta)(x)=\beta^\top h_t(x),
\]
with representation span $\mathcal S_t=\Ran(T_t)$. Let $C_t=T_t^\ast T_t=\mathbb E[h_th_t^\top]$ and $K_t=T_tT_t^\ast$. Their positive spectra coincide.

For a neuron subset $I\subset\{1,\ldots,m\}$ let $T_{t,I}$ keep only those coordinates and let $P_t,P_{t,I}$ denote the corresponding prediction-space projectors.

\begin{proposition}[Exact current cost of neuron contraction]\label{prop:neuron-cost}
The increase in population risk of the optimal frozen linear readout caused by retaining only $I$ is
\[
\boxed{
D_t(I)=\norm{(P_t-P_{t,I})f^\star}^2\ge0.}
\]
If $\Ran(T_{t,I})=\Ran(T_t)$, then $D_t(I)=0$.
\end{proposition}

\begin{proof}
The optimal frozen-readout excess risk in a representation span $\mathcal S$ is
$\norm{(I-P_\mathcal S)f^\star}^2$. Since $\Ran(T_{t,I})\subseteq\Ran(T_t)$, the projectors are nested and Pythagoras yields the difference $\norm{(P_t-P_{t,I})f^\star}^2$.
\end{proof}

\begin{theorem}[A rank-$r$ representation contains an exact $r$-neuron basis]\label{thm:subset-basis}
If $\mathrm{rank}(T_t)=r< m$, there exists a subset $I$ of exactly $r$ neuron coordinates such that
\[
\Ran(T_{t,I})=\Ran(T_t).
\]
Therefore the current optimal frozen-readout risk is preserved exactly after deleting $m-r$ neurons.
\end{theorem}

\begin{proof}
The $m$ coordinate feature functions $h_{t,1},\ldots,h_{t,m}$ span the $r$-dimensional space $\Ran(T_t)$. Every finite spanning set contains a basis subset of size $r$.
\end{proof}

\begin{remark}[Why not immediately compress to one feature?]
For a single scalar target, if the population-optimal readout coefficient were known, one could collapse the \emph{current predictor} into the scalar feature $\beta_t^{\star\top}h_t$. That does not preserve future feature-learning capacity. The relevant object for training-time contraction is therefore not current predictive sufficiency alone but \emph{current sufficiency plus a bound on future discovery value}. This is the central difference between deployment compression and training-time contraction.
\end{remark}

\section{The onset of safe reduction in the exact quadratic phase model}
The endogenous spectral-geometry monograph provides a solvable nonlinear teacher--student model in which the pruning-onset question can be answered cleanly. The learned covariance $A_t$ has a transported random bulk with upper edge $b_+(t)$ satisfying
\[
b_+(t)\sim(4t)^{-1},
\]
while every fixed teacher eigenvalue $\theta_q$ has a unique dynamic BBP-type crossing time $\tau_{\theta_q,\gamma}$. Before that time the teacher mode has no separated outlier; afterward a unique outlier $z_q(t)>b_+(t)$ appears, converges to $\theta_q$, and its eigenvector overlap $\omega_q(t)$ with the teacher direction converges to one. The tangent prediction operator sends this to a spike $\kappa_q(t)=4z_q(t)$ above the prediction bulk edge $4b_+(t)$.

\begin{theorem}[Oracle phase-aware pruning onset in the quadratic model]\label{thm:bbp-onset}
Consider finitely many distinct positive teacher modes $\theta_1,\ldots,\theta_k$ satisfying the supercritical terminal conditions of the quadratic phase theorem. Define
\[
\tau_{\rm sep}=\max_{1\le q\le k}\tau_{\theta_q,\gamma}.
\]
Then:
\begin{enumerate}[label=(\roman*)]
\item No rule that identifies all teacher modes solely as detached covariance/prediction outliers can be correct for every $t<\tau_{\rm sep}$.
\item For every $\varepsilon>0$, there exists $T_\varepsilon>\tau_{\rm sep}$ such that for all $t\ge T_\varepsilon$,
\[
b_+(t)\le\varepsilon,
\qquad
1-\omega_q(t)\le\varepsilon\quad\forall q.
\]
\item Consequently, for sufficiently large $t$ there exists a deterministic threshold $c_t$ that separates all teacher-associated prediction spikes from the random bulk:
\[
4b_+(t)<c_t<\min_q4z_q(t).
\]
\end{enumerate}
Thus the solvable model has an endogenous discovery phase followed by an asymptotically identifiable compression phase.
\end{theorem}

\begin{proof}
Part (i) follows because at least one teacher mode is, by definition of $\tau_{\rm sep}$, still absorbed in the bulk before the last crossing. For (ii), the phase theorem gives $b_+(t)\to0$ and $\omega_q(t)\to1$ for each fixed $q$; finiteness of $k$ lets one choose a common $T_\varepsilon$. It also gives $z_q(t)\to\theta_q>0$. Hence for sufficiently large $t$, $4b_+(t)<2\min_q\theta_q$ while $4z_q(t)>2\min_q\theta_q$ after increasing $T_\varepsilon$ if needed, so a separating threshold exists, proving (iii).
\end{proof}

\subsection{Why premature rank-one contraction creates a discovery delay}
The preceding onset theorem concerns observability. In the rank-one teacher model one can go further and quantify the optimization price of contracting before a target direction has macroscopic overlap.

\begin{proposition}[Exact target-angle dynamics after rank-one contraction]\label{prop:rankone-angle}
Let the teacher be $A^\star=\theta uu^\top$ with $\theta>0$, and after a contraction let the student have one active column $w_s$, so $A_s=w_sw_s^\top$. Define
\[
\omega_s=\frac{(u^\top w_s)^2}{\norm{w_s}^2}\in[0,1].
\]
Under the exact population gradient flow
\[
\dot w_s=-2(w_sw_s^\top-A^\star)w_s,
\]
for every interval on which $0<\omega_s<1$,
\[
\boxed{
\frac{1-\omega_{t+h}}{\omega_{t+h}}
=
\frac{1-\omega_t}{\omega_t}e^{-4\theta h}.}
\]
Therefore the additional time required to reach squared teacher overlap at least $1-\varepsilon$ is
\[
\boxed{
T_\varepsilon(\omega_t)
=
\frac{1}{4\theta}
\left[
\log\frac{(1-\omega_t)(1-\varepsilon)}{\omega_t\varepsilon}
\right]_+.}
\]
\end{proposition}

\begin{proof}
Decompose $w=cu+z$ with $z\perp u$ and write $q=\norm w^2$. The gradient flow gives
\[
\dot c=-2(q-\theta)c,
\qquad
\dot z=-2qz.
\]
Hence
\[
\frac{d}{ds}\log\frac{\norm z^2}{c^2}=-4\theta.
\]
Since $(1-\omega)/\omega=\norm z^2/c^2$, integration yields the odds identity. Solving the inequality $\omega_{t+h}\ge1-\varepsilon$ gives the hitting time.
\end{proof}

\begin{corollary}[Discovery-delay separation before and after spectral resolution]\label{cor:discovery-delay}
At isotropic initialization, any fixed sample covariance eigenvector is Haar relative to an independent deterministic teacher direction. Its squared overlap satisfies $\omega_0=O_P(d^{-1})$. Consequently a blind rank-one contraction at initialization requires
\[
T_\varepsilon(\omega_0)
=
\frac{\log d}{4\theta}+O_P(1)
\]
additional population time to obtain constant teacher overlap. By contrast, if contraction is performed at a fixed supercritical checkpoint $t>\tau_{\theta,\gamma}$ and retains the teacher-associated outlier, Theorem~5.5 of the endogenous phase theory gives $\omega_t\to\omega(t)>0$, so $T_\varepsilon(\omega_t)=O(1)$. As $t\to\infty$, $\omega(t)\to1$ and the additional alignment delay vanishes.
\end{corollary}

\begin{proof}
For a Haar unit vector $v$ independent of $u$, $d|u^\top v|^2$ is tight, so $\log(1/\omega_0)=\log d+O_P(1)$. Apply Proposition~\ref{prop:rankone-angle}. In the supercritical regime the outlier-overlap theorem supplies a nonzero deterministic limit, and its long-time limit is one.
\end{proof}

This result sharpens the meaning of ``do not prune too early.'' Premature contraction need not destroy the teacher forever: infinitesimal overlap can eventually be amplified. What it does is convert an order-one resolved feature into a dimension-dependent discovery delay. The BBP crossing is therefore a natural observability boundary for avoiding that delay, not a mystical universal threshold.

The theorem does \emph{not} claim that a BBP threshold is a universal pruning rule for arbitrary networks. It establishes the exact mechanism in one nonlinear model and provides the theory-derived reason for a warm-up/separation phase.

\section{Exact rank-one risk dynamics and the price of discovery}
The rank-one angle law identifies a dimension-dependent discovery delay, but by itself it does not say whether avoiding that delay is worth the compute of a wide exploratory phase. In the rank-one quadratic model the radial dynamics are also exactly solvable, so the discovery benefit can be converted into a critical wide-to-compact compute ratio without approximation.

Let $A^\star=\theta uu^\top$ with $\theta>0$, $\norm u=1$, and let the rank-one student be $A_s=w_sw_s^\top$. Define
\[
q_s:=\norm{w_s}^2,
\qquad
\omega_s:=\frac{(u^\top w_s)^2}{q_s},
\]
for $q_s>0$. The population risk is
\[
\R_1(q,\omega)
:=\frac12\norm{ww^\top-\theta uu^\top}_F^2
=\frac12\left(q^2+\theta^2-2\theta q\omega\right).
\]

\begin{proposition}[Closed-form rank-one state and risk]\label{prop:rankone-closed-form}
Under
\[
\dot w_s=-2(w_sw_s^\top-A^\star)w_s,
\]
with $q_0>0$ and $0<\omega_0<1$,
\[
\boxed{
\omega_s
=\frac{\omega_0e^{4\theta s}}
{1-\omega_0+\omega_0e^{4\theta s}},}
\]
and
\[
\boxed{
q_s
=
\frac{1-\omega_0+\omega_0e^{4\theta s}}
{q_0^{-1}+4(1-\omega_0)s+(\omega_0/\theta)(e^{4\theta s}-1)}.}
\]
Consequently the full rank-one risk path is explicit:
\[
\boxed{
\R_1(s;q_0,\omega_0)
=\frac12\left(q_s^2+\theta^2-2\theta q_s\omega_s\right).}
\]
Moreover $s\mapsto\R_1(s;q_0,\omega_0)$ is nonincreasing.
\end{proposition}

\begin{proof}
The odds identity in Proposition~\ref{prop:rankone-angle} gives the displayed logistic form for $\omega_s$. For the radial coordinate,
\[
\dot q_s
=2w_s^\top\dot w_s
=-4q_s(q_s-\theta\omega_s).
\]
Set $r_s=q_s^{-1}$. Then
\[
\dot r_s+4\theta\omega_s r_s=4.
\]
The logistic formula implies
\[
\exp\left(\int_0^s4\theta\omega_a\,da\right)
=1-\omega_0+\omega_0e^{4\theta s}.
\]
Multiplying the linear ODE for $r_s$ by this integrating factor and integrating gives
\[
r_s\bigl(1-\omega_0+\omega_0e^{4\theta s}\bigr)
=q_0^{-1}+4(1-\omega_0)s
+\frac{\omega_0}{\theta}(e^{4\theta s}-1),
\]
which yields the formula for $q_s$. The risk identity follows from
$\Tr((ww^\top)^2)=q^2$ and
$\Tr(ww^\top uu^\top)=q\omega$.
Finally the trajectory is gradient flow, hence
$\dot\R_1=-\norm{\nabla_w\R_1}^2\le0$.
\end{proof}

Now compare a wide discovery phase with a compact-from-start rank-one control. At wide time $t>0$, let a target-free structural rule contract the wide state to a rank-one state $(q_t^W,\omega_t^W)$. After an additional compact horizon $H\ge0$, define
\[
\bar R(t,H)
:=\R_1(H;q_t^W,\omega_t^W).
\]
Let the matched rank-one control start from $(q_0^C,\omega_0^C)$ and define its hitting time
\[
T_C(r)
:=\inf\{s\ge0:\R_1(s;q_0^C,\omega_0^C)\le r\}.
\]

\begin{theorem}[Critical compute ratio for temporary wide discovery]\label{thm:rankone-critical-compute}
Let one unit of wide training time cost $C_W$ and one unit of rank-one training time cost $C_1$, and write
\[
\rho:=\frac{C_W}{C_1}\ge1.
\]
The wide-then-rank-one path spends compute $C_Wt+C_1H$. At the same compute, the compact-from-start rank-one control can train for time
\[
H+\rho t.
\]
Therefore temporary wide discovery wins the equal-compute comparison exactly when
\[
\boxed{
\bar R(t,H)
<
\R_1(H+\rho t;q_0^C,\omega_0^C).}
\]
If the compact risk path is strictly decreasing until it first reaches $\bar R(t,H)$, this is equivalent to
\[
\boxed{
\rho<\rho_\star(t,H)
:=\frac{T_C(\bar R(t,H))-H}{t}.}
\]
Thus $\rho_\star$ is the maximum wide-to-compact cost ratio that the observed discovery advantage can repay. If $\rho_\star\le1$, temporary width cannot win even when a wide step costs no more than a compact step. If $\rho_\star>1$, the discovery mechanism has a nontrivial compute budget: it is beneficial on hardware/architectures whose realized cost ratio lies below $\rho_\star$ and not beneficial above it.
\end{theorem}

\begin{proof}
The compute equality
$C_Wt+C_1H=C_1(H+\rho t)$
is immediate from the definition of $\rho$. Hence the first boxed inequality is exactly the equal-compute risk comparison. Under strict monotonicity, the compact path has risk larger than $\bar R(t,H)$ precisely before time $T_C(\bar R(t,H))$. Therefore
\[
H+\rho t<T_C(\bar R(t,H))
\]
is necessary and sufficient for a strict wide-path win. Solving for $\rho$ yields the threshold.
\end{proof}

\begin{remark}[Why the critical ratio is the right controlled diagnostic]
A quadratic phase model may require many exploratory columns to produce a clean random-matrix bulk, while a real structured neural contraction may save a much smaller factor of compute. Hard-coding the number of exploratory columns as the compute ratio would conflate the statistical model used to generate spectral separation with a hardware cost model. The quantity $\rho_\star$ avoids that mistake. The experiment should report the discovery advantage and the entire critical-cost frontier first; only afterward should a measured architecture-specific $C_W/C_1$ be placed on that frontier.
\end{remark}

\section{A master finite-horizon safe-contraction theorem}
The previous results can be assembled into one theorem. This theorem is intentionally modular: it separates the value of the ideal prediction-space deletion, the motion of the learned geometry, and the error introduced when that ideal spectral decision is realized by actual structural coordinates.

Fix a contraction checkpoint $t$. Let $P_D(s)$ be a positive prediction-space spectral cluster of $L_s=J_sJ_s^\ast$ that is proposed for deletion, and let $Q_D(s)$ be its paired right-singular parameter projector. Put $Q_K(s)=I-Q_D(s)$. Let $S$ be a fixed structural projector selected at time $t$, representing the neurons, channels, heads, blocks, or rank directions that remain physically active after contraction. Define
\[
E_D(s)
:=\norm{Q_D(s)J_s^\ast e_s}^2
=\langle e_s,L_sP_D(s)e_s\rangle.
\]
At the contraction checkpoint define the structural mismatch
\[
\delta_0:=\norm{(I-S)Q_K(t)}_{\op}.
\]

\begin{theorem}[Master safe dynamic spectral reduction bound]\label{thm:master-safe}
Assume on $[t,t+H]$ that:
\begin{enumerate}[label=(\roman*)]
\item the full and structurally reduced trajectories remain in a common tube on which the gradient vector field $g(\theta)=J_\theta^\ast e_\theta$ is $L_g$-Lipschitz, $\norm{J_\theta}_{\op}\le B_J$, and both residual norms are at most $B_e$;
\item $\norm{D^2F(\theta_s)}\le B_H$ and $\norm{L_s}_{\op}\le B_L$;
\item the prediction cluster $P_D(s)$ remains separated from the rest of $L_s$ by at least $\Delta_L>0$;
\item the paired kept parameter cluster $Q_K(s)$ remains separated in $G_s=J_s^\ast J_s$ by at least $\Delta_G>0$.
\end{enumerate}
Let
\[
\Delta\R_{t,H}:=\R(\theta_t)-\R(\theta_{t+H})\le \R(\theta_t),
\]
\[
V_L(t,H):=2B_HB_J\sqrt{H\Delta\R_{t,H}},
\]
and
\[
\overline E_D(t,H)
:=E_D(t)
+B_e^2\left(1+\frac{2B_L}{\Delta_L}\right)V_L(t,H).
\]
Furthermore define
\[
\overline\delta_S(t,H)
:=\delta_0+
\frac{4B_HB_J}{\Delta_G}\sqrt{H\Delta\R_{t,H}}.
\]
If the post-contraction initial embedding error is $d_0=\norm{\theta_t-\widetilde\theta_t}$, then
\[
\boxed{
\begin{aligned}
\norm{\theta_{t+H}-\widetilde\theta_{t+H}}
\le e^{L_gH}\Big[&d_0
+H\sqrt{\overline E_D(t,H)}\\
&+\overline\delta_S(t,H)\sqrt{H\Delta\R_{t,H}}
\Big].
\end{aligned}}
\]
Consequently,
\[
\boxed{
\begin{aligned}
|\R(\theta_{t+H})-\R(\widetilde\theta_{t+H})|
\le B_eB_Je^{L_gH}\Big[&d_0
+H\sqrt{\overline E_D(t,H)}\\
&+\overline\delta_S(t,H)\sqrt{H\Delta\R_{t,H}}
\Big].
\end{aligned}}
\]
The same bracket multiplied by $B_J$ bounds the terminal prediction discrepancy.
\end{theorem}

\begin{proof}
At any $s\in[t,t+H]$, structural transfer gives
\[
\norm{(I-S)g(\theta_s)}
\le
\norm{Q_D(s)g(\theta_s)}
+\norm{(I-S)Q_K(s)}_{\op}\norm{g(\theta_s)}.
\]
The first term is $\sqrt{E_D(s)}$. The spectral-persistence theorem yields
\[
E_D(s)\le\overline E_D(t,H).
\]
The fixed-structure persistence theorem gives
\[
\norm{(I-S)Q_K(s)}_{\op}
\le\overline\delta_S(t,H).
\]
Therefore
\[
\int_t^{t+H}\norm{(I-S)g(\theta_s)}\,ds
\le
H\sqrt{\overline E_D(t,H)}
+\overline\delta_S(t,H)
\int_t^{t+H}\norm{g(\theta_s)}\,ds.
\]
By exact loss dissipation,
\[
\int_t^{t+H}\norm{g(\theta_s)}^2ds
=\Delta\R_{t,H},
\]
so Cauchy--Schwarz gives
\[
\int_t^{t+H}\norm{g(\theta_s)}ds
\le\sqrt{H\Delta\R_{t,H}}.
\]
Finally, apply the nonlinear structural-flow Gr\"onwall theorem. Since
$e^{L_g(t+H-s)}\le e^{L_gH}$, its weighted omitted-gradient integral is at most $e^{L_gH}$ times the preceding unweighted bound. The terminal prediction and risk inequalities then follow from the $B_J$-Lipschitz predictor bound and the residual bound exactly as in Theorem~\ref{thm:gronwall}.
\end{proof}

\begin{corollary}[Asymptotically safe compression regime]\label{cor:asymptotic-safe}
Consider a sequence of contraction checkpoints for which, for a fixed horizon $H$ and uniformly bounded constants in Theorem~\ref{thm:master-safe},
\[
d_0\to0,
\qquad
E_D(t)\to0,
\qquad
\delta_0\to0,
\qquad
\Delta\R_{t,H}\to0.
\]
Then the terminal prediction and risk discrepancy between dense and structurally reduced training converges to zero.
\end{corollary}

\begin{proof}
Every term in the bracket of Theorem~\ref{thm:master-safe} converges to zero under the stated uniform bounds.
\end{proof}

\begin{remark}[What the theorem does and does not certify]
The theorem is not an eigenvalue threshold. It requires four distinct pieces of evidence: small target-weighted energy in the sector proposed for deletion, little reachable spectral motion over the decision horizon, a persistent gap that makes the sector observable, and a small structural-realization error. This is precisely why a small weight singular value, a small prediction eigenvalue, or a stable pruning mask alone cannot constitute the general certificate.
\end{remark}

\begin{remark}[Observable versus oracle quantities]
The theorem is a population statement. In practice $E_D(t)$, gaps, target mass, and structural mismatch must be estimated. A method becomes certified only after adding finite-sample confidence bounds for these quantities. Plugging point estimates or extrapolated recent velocities into the theorem produces a useful controller proxy, not a theorem-level certificate.
\end{remark}

\section{A checkpoint-computable prediction-gap certificate}
The master theorem is still partly oracle because it is written with the unknown future loss drop $\Delta\R_{t,H}$ and assumes that the relevant spectral gaps remain open over the horizon. Both quantities can be removed from the decision rule. The key observation is that square-loss dissipation already controls how far the dense parameter trajectory can travel after the checkpoint.

For the dense flow put
\[
g_s=J_s^\ast e_s,
\qquad
A_t(h):=\int_t^{t+h}\norm{g_s}\,ds.
\]
Let $L_g$ be a valid Lipschitz constant of the gradient vector field on the declared local tube. Then two current-state bounds are available:
\[
A_t(H)
\le \sqrt{H\R(\theta_t)},
\]
by Cauchy--Schwarz and exact loss dissipation, and
\[
A_t(H)
\le \norm{g_t}\,\phi(L_g,H),
\qquad
\phi(L,H):=
\begin{cases}
(e^{LH}-1)/L,&L>0,\\
H,&L=0,
\end{cases}
\]
because $\norm{g_{t+h}}\le e^{L_gh}\norm{g_t}$. Define the checkpoint path budget
\[
\boxed{
\overline A_t(H)
:=
\min\left\{
\sqrt{H\R(\theta_t)},
\norm{g_t}\phi(L_g,H)
\right\}.}
\]
Thus late in training the current gradient norm can make the horizon budget much smaller than the crude loss-only envelope.

\begin{theorem}[Checkpoint-computable horizon certificate]\label{thm:checkpoint-certificate}
Fix a checkpoint $t$, horizon $H>0$, a prediction-space cluster $P_D(t)$ proposed for deletion, its paired kept parameter cluster $Q_K(t)$, and a fixed structural projector $S$. Let
\[
E_D(t)=\norm{Q_D(t)g_t}^2,
\qquad
\delta_0=\norm{(I-S)Q_K(t)}_{\op},
\qquad
 d_0=\norm{\theta_t-\widetilde\theta_t}.
\]
Assume:
\begin{enumerate}[label=(\roman*)]
\item on the dense horizon, $\norm{D^2F}\le B_H$;
\item on a common tube containing the dense and structurally reduced continuations, $g$ is $L_g$-Lipschitz and $F$ is $B_{J,\mathrm{tube}}$-Lipschitz;
\item at the checkpoint the discarded prediction cluster and paired kept parameter cluster have gaps $\Delta_{L,0}>0$ and $\Delta_{G,0}>0$.
\end{enumerate}
Let
\[
A:=\overline A_t(H),
\qquad
\overline B_J:=\norm{J_t}_{\op}+B_HA,
\qquad
V:=2B_H\overline B_JA.
\]
Define the a priori gap floors
\[
\underline\Delta_L:=\Delta_{L,0}-2V,
\qquad
\underline\Delta_G:=\Delta_{G,0}-2V.
\]
If
\[
\boxed{
\underline\Delta_L>0,
\qquad
\underline\Delta_G>0,}
\]
then both clusters remain isolated throughout $[t,t+H]$. Moreover, with
\[
B_e:=\sqrt{2\R(\theta_t)},
\qquad
\overline B_L:=\overline B_J^2,
\]
set
\[
\overline E_D
:=E_D(t)
+B_e^2
\left(1+\frac{2\overline B_L}{\underline\Delta_L}\right)V.
\]
The dense omitted-gradient path of the structural action satisfies
\[
\boxed{
\int_t^{t+H}\norm{(I-S)g_s}\,ds
\le
H\sqrt{\overline E_D}
+\delta_0A
+\frac{2B_H\overline B_J}{\underline\Delta_G}A^2.}
\]
Consequently the terminal prediction discrepancy obeys the fully checkpoint-computable bound
\[
\boxed{
\begin{aligned}
\norm{F(\theta_{t+H})-F(\widetilde\theta_{t+H})}
\le
\mathfrak D^{\rm chk}_{t,H}(S)
:=B_{J,\mathrm{tube}}e^{L_gH}
\Bigg[
&d_0+H\sqrt{\overline E_D}
+\delta_0A\\
&+\frac{2B_H\overline B_J}{\underline\Delta_G}A^2
\Bigg].
\end{aligned}}
\]
Every quantity on the right is fixed or measured at the checkpoint, apart from the declared deterministic tube/curvature bounds.
\end{theorem}

\begin{proof}
Exact dissipation gives
\[
\int_t^{t+H}\norm{g_s}^2ds
=\R(\theta_t)-\R(\theta_{t+H})
\le\R(\theta_t),
\]
so Cauchy--Schwarz yields the first path bound. Since $g$ is $L_g$-Lipschitz along its own flow,
\[
\frac{d}{dh}\norm{g_{t+h}}
\le L_g\norm{g_{t+h}}
\]
almost everywhere, which gives the exponential current-gradient bound and therefore $A_t(H)\le A$.

Along the dense path,
\[
\norm{J_s-J_t}_{\op}
\le
B_H\int_t^s\norm{g_r}\,dr
\le B_HA,
\]
so $\norm{J_s}_{\op}\le\overline B_J$. Therefore both
$L_s=J_sJ_s^\ast$ and $G_s=J_s^\ast J_s$ satisfy
\[
\int_t^{t+H}\norm{\dot L_s}_{\op}ds
\le V,
\qquad
\int_t^{t+H}\norm{\dot G_s}_{\op}ds
\le V.
\]
Weyl's perturbation inequality moves each side of an inter-cluster gap by at most the operator perturbation, so each gap can shrink by at most $2V$. Positivity of the displayed gap floors therefore certifies persistence rather than assuming it.

Because the dense residual norm is nonincreasing, $\norm{e_s}\le B_e$, while
$\norm{L_s}_{\op}\le\overline B_J^2=\overline B_L$. The discarded-energy persistence theorem with the certified gap floor then gives
\[
E_D(s)\le\overline E_D.
\]
For the parameter projector, let
\[
A(s):=\int_t^s\norm{g_r}\,dr.
\]
The Riesz-projector bound and the same curvature calculation imply
\[
\norm{Q_K(s)-Q_K(t)}_{\op}
\le
\frac{4B_H\overline B_J}{\underline\Delta_G}A(s).
\]
Structural transfer therefore yields
\[
\norm{(I-S)g_s}
\le
\sqrt{\overline E_D}
+\left[
\delta_0+
\frac{4B_H\overline B_J}{\underline\Delta_G}A(s)
\right]\norm{g_s}.
\]
Integrating and using $A'(s)=\norm{g_s}$ almost everywhere gives
\[
\int_t^{t+H}\norm{(I-S)g_s}ds
\le
H\sqrt{\overline E_D}
+\delta_0A
+\frac{4B_H\overline B_J}{\underline\Delta_G}
\int_t^{t+H}A(s)A'(s)ds,
\]
which is exactly the omitted-gradient bound because
$\int A A'=A(H)^2/2\le A^2/2$.
Finally apply Theorem~\ref{thm:gronwall} and upper-bound its exponential kernel by $e^{L_gH}$. The common-tube predictor Lipschitz constant gives the terminal prediction bound.
\end{proof}

\begin{remark}[What is new relative to the master theorem]
The master theorem is structurally correct but asks for two future facts: $\Delta\R_{t,H}$ and persistence of the spectral gaps. Theorem~\ref{thm:checkpoint-certificate} replaces both by a current-state path budget. The practical gate is now explicit:
\[
\boxed{
2V<\min\{\Delta_{L,0},\Delta_{G,0}\}
}
\]
before any projector-based deletion argument is allowed to proceed.
\end{remark}

\begin{remark}[Why the target-weighted term remains indispensable]
Even after the gap is certified, the leading contribution is $H\sqrt{\overline E_D}$. Hence a spectrally stable low-eigenvalue sector is not enough: the sector must also carry little target-weighted gradient energy. Conversely, if $E_D(t)$ is small but the curvature path budget is large, $\overline E_D$ warns that future feature motion can still make the sector valuable.
\end{remark}

\begin{corollary}[Checkpoint-to-observability bridge]\label{cor:checkpoint-observability}
For clipped squared loss $\ell_B$, let
\[
M_{2}^{\rm chk}(S)
:=
\min\left\{
B^2,
4B\left[\mathfrak D^{\rm chk}_{t,H}(S)\right]^2
\right\}.
\]
Then $M_2^{\rm chk}(S)$ is a deterministic pre-probe second-moment bound for the paired structural loss difference. It can therefore be inserted directly into Theorem~\ref{thm:localized-damage}. The resulting UCB is a finite-sample certificate driven by the current target-weighted spectrum rather than by the global loss range alone.
\end{corollary}

\begin{proof}
Theorem~\ref{thm:checkpoint-certificate} supplies the population $L^2$ prediction-gap bound. Lemma~\ref{lem:loss-gap-moment} converts it into the paired-loss second-moment bound, and Theorem~\ref{thm:localized-damage} applies because the geometric quantity is fixed before observing the independent probe.
\end{proof}

\begin{remark}[Remaining non-oracle constant]
The only genuinely architecture-dependent analytic burden left in this certificate is to obtain useful deterministic local bounds on $B_H$, $L_g$, and $B_{J,\mathrm{tube}}$. Estimating them by a point sample is a controller proxy; certifying them requires an architecture-specific norm, interval, PAC-Bayes, Lipschitz, or other valid local bound. This is now a sharply isolated problem rather than an unspecified ``future motion'' term.
\end{remark}

\section{When the spectral bulk is ready to be removed}
The general checkpoint certificate treats an arbitrary discarded spectral cluster.  For the central use case of this paper we can say more.  The proposed deleted set is the lower spectral bulk, while the resolved outliers are kept.  This structure gives a second, target-mass-based envelope on the future value of the bulk and makes the question ``when should the bulk be removed?'' mathematically explicit.

Fix a checkpoint $t$.  Suppose the spectrum of the positive prediction operator $L_t$ is split into a lower cluster
\[
\sigma_B(L_t)\subset[0,b_t]
\]
and a resolved upper cluster
\[
\sigma_S(L_t)\subset[s_t,\infty),
\qquad
\Delta_t:=s_t-b_t>0.
\]
Let $P_B(t)$ be the spectral projector of the lower cluster and let $P_B(s)$ denote its Riesz continuation as long as the gap stays open.  Define the residual target mass and gradient energy in the bulk by
\[
M_B(s):=\norm{P_B(s)e_s}^2,
\qquad
E_B(s):=\ip{e_s}{L_sP_B(s)e_s}.
\]
The distinction is important: $M_B$ asks how much unresolved target residual lies in the bulk, whereas $E_B$ asks how much instantaneous loss dissipation that mass can buy through the current eigenvalues.

\begin{theorem}[Target-mass persistence of an isolated lower bulk]\label{thm:bulk-mass}
Assume on $[t,t+H]$ that $L_s$ is differentiable and positive semidefinite, $\norm{e_s}\le B_e$, and
\[
\Gamma_{t,H}:=\int_t^{t+H}\norm{\dot L_s}_{\op}\,ds
<\frac{\Delta_t}{2}.
\]
Then throughout the horizon:
\begin{enumerate}[label=(\roman*)]
\item the lower and upper clusters remain separated by at least
\[
\underline\Delta_t(H):=\Delta_t-2\Gamma_{t,H}>0;
\]
\item the upper edge of the continued bulk obeys
\[
\lambda_{\max}(L_s|_{\Ran P_B(s)})\le b_t+\Gamma_{t,H};
\]
\item the bulk projector path length obeys
\[
\int_t^{t+H}\norm{\dot P_B(s)}_{\op}\,ds
\le
\frac{2\Gamma_{t,H}}{\underline\Delta_t(H)};
\]
\item the target mass that can enter the bulk is bounded by
\[
\boxed{
M_B(s)
\le
\overline M_B(t,H)
:=M_B(t)+
\frac{2B_e^2\Gamma_{t,H}}{\underline\Delta_t(H)};}
\]
\item consequently the entire future bulk gradient energy satisfies
\[
\boxed{
E_B(s)
\le
\overline E_B^{\rm mass}(t,H)
:=(b_t+\Gamma_{t,H})\overline M_B(t,H).}
\]
In particular,
\[
\boxed{
\int_t^{t+H}E_B(s)\,ds
\le H\overline E_B^{\rm mass}(t,H),
\qquad
\int_t^{t+H}\sqrt{E_B(s)}\,ds
\le H\sqrt{\overline E_B^{\rm mass}(t,H)}.}
\]
\end{theorem}

\begin{proof}
For every $h\le H$,
\[
\norm{L_{t+h}-L_t}_{\op}
\le
\int_t^{t+h}\norm{\dot L_s}_{\op}ds
\le\Gamma_{t,H}.
\]
Weyl's perturbation inequality therefore moves the upper edge of the lower cluster upward by at most $\Gamma_{t,H}$ and the lower edge of the upper cluster downward by at most the same amount.  This proves (i)--(ii).  Along the resulting isolated Riesz cluster,
\[
\norm{\dot P_B(s)}_{\op}
\le
\frac{2\norm{\dot L_s}_{\op}}{\underline\Delta_t(H)},
\]
which integrates to (iii).

Now differentiate $M_B=\ip e{P_Be}$.  Since $\dot e=-Le$ and $P_B$ commutes with $L$,
\[
\dot M_B
=-2\ip{P_Be}{LP_Be}+\ip e{\dot P_Be}
\le
\norm e^2\norm{\dot P_B}_{\op}.
\]
The first term is nonpositive because $L\succeq0$.  Integrating, using $\norm e\le B_e$, and applying (iii) gives (iv).  Finally,
\[
E_B(s)=\ip{P_Be_s}{L_sP_Be_s}
\le
\lambda_{\max}(L_s|_{\Ran P_B(s)})M_B(s),
\]
so (ii) and (iv) imply (v).  The two integral bounds follow immediately from the uniform pointwise envelope.
\end{proof}

\begin{remark}[The bulk is a discovery reservoir before it is nuisance]
The theorem gives a precise version of the central intuition.  A small bulk edge $b_t$ is not enough for deletion.  Early in training, $\Gamma_{t,H}$ can be large, the gap may not be persistent, and a substantial amount of target mass can still rotate into the lower cluster.  Only after the outliers are separated and the reachable projector path becomes small does the same lower bulk acquire a uniformly small future gradient-energy budget.
\end{remark}

The target-mass envelope complements the discarded-energy envelope in Theorem~\ref{thm:checkpoint-certificate}.  For a lower bulk we may use whichever valid envelope is tighter.  Using the notation of that theorem, let
\[
V=2B_H\overline B_JA,
\qquad
\underline\Delta_L=\Delta_{L,0}-2V,
\]
and suppose $\underline\Delta_L>0$.  Since the dense residual norm is bounded by $B_e=\sqrt{2\R(\theta_t)}$, define
\[
\overline M_B^{\rm chk}
:=M_B(t)+\frac{2B_e^2V}{\underline\Delta_L},
\]
\[
\overline E_B^{\rm mass,chk}
:=(b_t+V)\overline M_B^{\rm chk}.
\]
The existing energy-persistence argument gives
\[
\overline E_B^{\rm energy,chk}
:=E_B(t)+B_e^2
\left(1+\frac{2\overline B_L}{\underline\Delta_L}\right)V.
\]
Both are valid uniform upper bounds on $E_B(s)$, hence so is
\[
\boxed{
\overline E_B^{\rm chk}
:=\min\left\{
\overline E_B^{\rm energy,chk},
\overline E_B^{\rm mass,chk}
\right\}.}
\]

\begin{corollary}[Checkpoint bulk-readiness certificate]\label{cor:bulk-readiness}
Under the assumptions of Theorem~\ref{thm:checkpoint-certificate}, specialize the discarded prediction cluster to the isolated lower bulk above and replace $\overline E_D$ there by $\overline E_B^{\rm chk}$.  Then the omitted-gradient path obeys the sharper bound
\[
\boxed{
\mathfrak G_B(t,H;S)
:=H\sqrt{\overline E_B^{\rm chk}}
+\delta_0A
+\frac{2B_H\overline B_J}{\underline\Delta_G}A^2,}
\]
and the terminal prediction discrepancy satisfies
\[
\boxed{
\mathfrak D_B(t,H;S)
:=B_{J,\mathrm{tube}}e^{L_gH}
\left[d_0+\mathfrak G_B(t,H;S)\right].}
\]
Thus a lower spectral bulk is geometrically ready for structural deletion over horizon $H$ only when the gap-persistence gate is open and $\mathfrak D_B(t,H;S)$ is inside the predeclared prediction-damage budget.
\end{corollary}

\begin{proof}
The proof of Theorem~\ref{thm:checkpoint-certificate} uses $\overline E_D$ only as a uniform upper bound on $E_D(s)$.  Theorem~\ref{thm:bulk-mass} supplies a second valid uniform bound for a lower cluster, so their minimum may be substituted without changing any subsequent step.  Structural transfer and Gr\"onwall then give the displayed quantities.
\end{proof}

\begin{definition}[Geometric bulk-onset time]\label{def:bulk-onset}
Fix a horizon $H$, a family of realizable structural actions $S_t$, and tolerances $\varepsilon_{\rm value},\varepsilon_{\rm pred}>0$.  The geometric onset time is
\[
\boxed{
\tau_{\rm bulk}
:=\inf\left\{t:\begin{array}{l}
\underline\Delta_L(t,H)>0,\quad \underline\Delta_G(t,H)>0,\\[2pt]
H\overline E_B^{\rm chk}(t,H)\le\varepsilon_{\rm value},\\[2pt]
\mathfrak D_B(t,H;S_t)\le\varepsilon_{\rm pred}
\end{array}\right\}.}
\]
This is a population geometric onset, not yet the final algorithmic trigger.  The irreversible contraction is executed only after the finite-sample observability and compute-value gates also pass.
\end{definition}

\begin{remark}[What the onset means]
The first line says that ``bulk'' and ``outliers'' are predicted to remain identifiable over the decision horizon.  The second says that, even along the dense trajectory, the lower cluster can buy only a small amount of future loss decrease.  The third says that actually removing the corresponding structural capacity cannot move the future predictor by more than the declared amount.  The later finite-sample UCB then asks whether this small predicted damage is observable tightly enough to trade it against real compute savings.
\end{remark}

\begin{remark}[Why this is not an eigenvalue threshold]
If $b_t$ is tiny but $M_B(t)$ is large, the bulk can still matter.  If both are tiny but $V$ is large, hidden signal can still rotate into it.  If all prediction-space quantities are favorable but the structural mismatch $\delta_0$ or contraction jump $d_0$ is large, the physical deletion can still be unsafe.  The rule is therefore
\[
\boxed{
\text{resolved bulk}
+\text{small target mass/value}
+\text{small reachable motion}
+\text{small structural mismatch}
\Longrightarrow
\text{ready to test for deletion}.}
\]

\section{An exact structural bridge: annealable rank-factor layers}
The general spectral theory still has a structural-realizability problem: a prediction eigenvector is not a literal neuron or matrix entry. A useful architecture should make an irreversible rank contraction an exact coordinate operation rather than an approximate analogy. One simple construction is an additive rank-factor layer.

Consider an affine layer
\[
z=W h+b,
\qquad
W=\sum_{q=1}^{r}s_q u_qv_q^\top,
\]
where $u_q\in\mathbb R^{d_{\rm out}}$, $v_q\in\mathbb R^{d_{\rm in}}$, and $s_q\in\mathbb R$ are trainable. No orthogonality is assumed here; the $q$'s are structural rank-one components, not automatically singular vectors. Let the remaining network parameters be denoted by $\xi$.

The raw factorization has a gauge symmetry that must be fixed before scale-based structural diagnostics are interpreted. For every nonzero $c,d$,
\[
(s_q,u_q,v_q)
\mapsto
\left(\frac{s_q}{cd},cu_q,dv_q\right)
\]
leaves $s_qu_qv_q^\top$ unchanged. Hence raw $|s_q|$ and Euclidean parameter-gradient norms are not representation invariant.

\begin{proposition}[Canonical component gauge]\label{prop:rank-gauge}
Every nondegenerate rank-one component admits the function-preserving canonicalization
\[
\bar u_q=\frac{u_q}{\norm{u_q}},
\qquad
\bar v_q=\frac{v_q}{\norm{v_q}},
\qquad
\bar s_q=s_q\norm{u_q}\norm{v_q}.
\]
Then
\[
\norm{\bar u_q}=\norm{\bar v_q}=1,
\qquad
|\bar s_q|=\norm{s_qu_qv_q^\top}_{F}.
\]
Thus the magnitude of the canonical scale is the gauge-invariant Frobenius magnitude of the represented component. Canonicalization changes neither the layer matrix nor the network predictor.
\end{proposition}

\begin{proof}
Substitution gives $\bar s_q\bar u_q\bar v_q^\top=s_qu_qv_q^\top$. The Frobenius norm of an outer product is the product of the vector norms, which gives the second identity.
\end{proof}

All scale and group-gradient diagnostics below are understood in this canonical gauge. A practical implementation can canonicalize a copied checkpoint before evaluating a candidate and reset the optimizer in both matched dense and reduced counterfactual continuations. This avoids attributing scientific meaning to an arbitrary factor scaling.

For a kept set $K\subset\{1,\ldots,r\}$ and dropped set $D=K^c$, define the physical compact layer
\[
W_K=\sum_{q\in K}s_qu_qv_q^\top.
\]
In the full parameter space define the \emph{ghost embedding} by leaving $u_q,v_q$ unchanged for $q\in D$ but setting $s_q=0$. Let $S_K$ be the coordinate projector that keeps $\xi,b$ and all parameters $(s_q,u_q,v_q)$ for $q\in K$ while freezing every parameter of the dropped groups.

\begin{theorem}[Exact physical embedding of rank contraction]\label{thm:rank-factor-embedding}
For every checkpoint and every kept set $K$:
\begin{enumerate}[label=(\roman*)]
\item the predictor of the physically compact rank-$|K|$ layer is exactly equal to the predictor of its ghost embedding in the rank-$r$ parameter space;
\item the projected gradient flow of the ghost model under $S_K$ is exactly the embedded gradient flow of the physically compact model;
\item if the compact initialization copies every kept/shared parameter and only zeros the dropped scales, then the full-space contraction jump is
\[
\boxed{
d_0=\norm{s_D}_2;}
\]
in the chosen gauge; in the canonical gauge each $|s_q|$ is also the component Frobenius magnitude;
\item deleting one component removes exactly $d_{\rm in}+d_{\rm out}+1$ trainable factor parameters and removes its two rank-factor matrix-vector contractions from the physical forward/backward graph.
\end{enumerate}
\end{theorem}

\begin{proof}
When $s_q=0$, the rank-one matrix $s_qu_qv_q^\top$ is identically zero, so removing the corresponding term or retaining it as a zero-scale ghost gives the same affine map and therefore the same network predictor. Derivatives with respect to every shared and kept parameter are also identical in the two representations. The projected ghost flow freezes the deleted coordinates, including the scale directions that could otherwise regrow, so its active-coordinate ODE is exactly the compact-model ODE. The only coordinates changed at the embedding checkpoint are $s_q$, $q\in D$, giving the Euclidean distance $\norm{s_D}_2$. The parameter-count statement follows directly from the shapes of $u_q,s_q,v_q$.
\end{proof}

This construction solves the structural bridge for one concrete family: rank can be reduced during training with a literal decrease in parameters and arithmetic. It does \emph{not} say that small $|s_q|$ alone makes a component useless. Target value is still required.

At a canonicalized checkpoint define the exact structural group-gradient energy
\[
E_D^{\rm grp}(t)
:=\norm{(I-S_K)g_t}^2
=
\sum_{q\in D}
\left[
\norm{\nabla_{u_q}\R_t}^2
+(\partial_{s_q}\R_t)^2
+\norm{\nabla_{v_q}\R_t}^2
\right].
\]
Because the groups occupy orthogonal parameter coordinates, this identity is exact in the fixed canonical coordinate representation. It is a structural analogue of the target-weighted modal energy $\mu_ja_j^2$: it measures how much present loss-dissipation budget actually flows through the parameters that would disappear.

\begin{theorem}[Gap-free checkpoint certificate for physical parameter groups]\label{thm:group-certificate}
Let $S$ be any fixed coordinate/group projector representing a physically realizable contraction, and suppose on a common tube containing the dense and projected trajectories that $g$ is $L_g$-Lipschitz and $F$ is $B_{J,\mathrm{tube}}$-Lipschitz. Define
\[
E_S(t):=\norm{(I-S)g_t}^2,
\qquad
A_t(h):=\int_t^{t+h}\norm{g_s}\,ds.
\]
As in Theorem~\ref{thm:checkpoint-certificate}, set
\[
\overline A_t(H)
=\min\left\{
\sqrt{H\R_t},
\norm{g_t}\phi(L_g,H)
\right\}.
\]
For the integral of the path budget define
\[
\overline B_t(H)
:=
\min\left\{
\frac23H^{3/2}\sqrt{\R_t},
\norm{g_t}\psi(L_g,H)
\right\},
\]
where
\[
\psi(L,H)
:=
\begin{cases}
\dfrac{e^{LH}-1-LH}{L^2},&L>0,\\[6pt]
H^2/2,&L=0.
\end{cases}
\]
Then the omitted dense gradient of the physical groups satisfies
\[
\boxed{
\int_t^{t+H}\norm{(I-S)g_s}\,ds
\le
H\sqrt{E_S(t)}+L_g\overline B_t(H).}
\]
If the contraction has initial full-space embedding error $d_0$, the terminal prediction discrepancy satisfies
\[
\boxed{
\norm{F(\theta_{t+H})-F(\widetilde\theta_{t+H})}
\le
B_{J,\mathrm{tube}}e^{L_gH}
\left[
d_0+H\sqrt{E_S(t)}+L_g\overline B_t(H)
\right].}
\]
For the canonically gauged rank-factor contraction of Theorem~\ref{thm:rank-factor-embedding}, one may take $d_0=\norm{s_D}_2$ and $E_S(t)=E_D^{\rm grp}(t)$ exactly.
\end{theorem}

\begin{proof}
Because $I-S$ is an orthogonal projector and $g$ is Lipschitz,
\[
\begin{aligned}
\norm{(I-S)g_s}
&\le
\norm{(I-S)g_t}
+\norm{(I-S)(g_s-g_t)}\\
&\le
\sqrt{E_S(t)}+L_g\norm{\theta_s-\theta_t}\\
&\le
\sqrt{E_S(t)}+L_gA_t(s-t).
\end{aligned}
\]
Integrating gives
\[
\int_t^{t+H}\norm{(I-S)g_s}ds
\le H\sqrt{E_S(t)}+L_g\int_0^H A_t(h)dh.
\]
The loss-dissipation path bound gives $A_t(h)\le\sqrt{h\R_t}$, whose integral is $\frac23H^{3/2}\sqrt{\R_t}$. The current-gradient Lipschitz bound gives
$A_t(h)\le\norm{g_t}\phi(L_g,h)$, whose integral is $\norm{g_t}\psi(L_g,H)$. Taking the smaller valid integral envelope gives $\overline B_t(H)$. The prediction statement is Theorem~\ref{thm:gronwall} followed by the common-tube predictor Lipschitz bound. The rank-factor specialization follows from Theorem~\ref{thm:rank-factor-embedding} and the orthogonal group-energy identity.
\end{proof}

\begin{remark}[Why this is useful]
The spectral checkpoint theorem and the structural-group theorem play different roles. The spectral theorem answers \emph{which capacity sector has little target-weighted future value and is unlikely to rotate back into relevance}. The group theorem answers \emph{whether the actual parameters we intend to delete are currently dynamically inactive enough to remove}. A practical rank-annealing controller should require both whenever both are available. This avoids the unjustified step ``small prediction eigenvalue $\Rightarrow$ delete arbitrary weights.''
\end{remark}

\begin{remark}[The scale is not a saliency score]
Even in canonical gauge, the term $\norm{s_D}_2$ controls the instantaneous embedding jump; it does not measure future usefulness. A component with tiny scale but large $|\partial_{s_q}\R|$ is exactly a component poised to regrow and should fail the group-gradient gate. Conversely a moderate scale can be removable after a function-preserving reparameterization or readout redistribution, but that requires a different embedding argument. The controller should therefore monitor canonical component magnitude, target-weighted group gradient, prediction-space discovery state, and finite-horizon risk jointly.
\end{remark}

\section{Finite-sample observability of structural damage}
The population master theorem is useful only if the risk price of a structural action can be resolved from data. The first counterfactual audit shows that this is not automatic: when the structural effect is small, two finite evaluation samples can even disagree on its sign. The correct finite-sample object is therefore the \emph{paired} loss difference, not two separately estimated risks.

Fix a checkpoint and a finite candidate family of structural actions $S_1,\ldots,S_M$. For candidate $m$, let $f_m$ denote the predictor obtained after the declared contraction/readout/training protocol over horizon $H$, and let $f_0$ denote the matched dense continuation. Assume these predictors are constructed without using an independent probe $\mathcal P=\{Z_i\}_{i=1}^n$. Let $\ell_B\in[0,B]$ be a bounded or clipped evaluation loss and define
\[
Z_{i,m}:=\ell_B(f_m,Z_i)-\ell_B(f_0,Z_i)\in[-B,B].
\]
The population structural damage and its paired probe estimate are
\[
D_m:=\mathbb E Z_{i,m},
\qquad
\widehat D_m:=\frac1n\sum_{i=1}^n Z_{i,m}.
\]
Let
\[
\widehat v_m:=\frac1{n-1}\sum_{i=1}^n(Z_{i,m}-\widehat D_m)^2
\]
be the empirical variance of the paired differences.

\begin{theorem}[Uniform paired empirical-Bernstein structural-damage bound]\label{thm:paired-damage}
Suppose $n\ge2$. With probability at least $1-\delta$ over the independent probe, simultaneously for all $m=1,\ldots,M$,
\[
\boxed{
D_m
\le
\widehat D_m
+
\sqrt{\frac{2\widehat v_m\log(2M/\delta)}{n}}
+
\frac{14B\log(2M/\delta)}{3(n-1)}.
}
\]
The same bound applied to $-Z_{i,m}$ gives a simultaneous two-sided confidence interval with the corresponding union-bound adjustment.
\end{theorem}

\begin{proof}
For each $m$, transform the paired variable to
\[
Y_{i,m}:=\frac{Z_{i,m}+B}{2B}\in[0,1].
\]
Its empirical variance is $\widehat v_m/(4B^2)$. The empirical Bernstein inequality gives, with failure probability $\delta/M$,
\[
\mathbb E Y_m
\le
\overline Y_m
+
\sqrt{\frac{2\widehat{\mathrm{Var}}(Y_m)\log(2M/\delta)}{n}}
+
\frac{7\log(2M/\delta)}{3(n-1)}.
\]
Multiplying by $2B$ and subtracting $B$ yields the displayed inequality. A union bound over $m$ completes the proof.
\end{proof}

\begin{corollary}[Finite-sample safe structural action]\label{cor:paired-safe}
Let $b_m\ge0$ be a predeclared admissible population risk budget for candidate $m$, possibly equal to its risk-equivalent compute value. Define
\[
U_m
:=
\widehat D_m
+
\sqrt{\frac{2\widehat v_m\log(2M/\delta)}{n}}
+
\frac{14B\log(2M/\delta)}{3(n-1)}.
\]
On the event of Theorem~\ref{thm:paired-damage}, every candidate satisfying
\[
\boxed{U_m\le b_m}
\]
has population structural damage $D_m\le b_m$. In particular, if $b_m=\tau\Delta C_m$ is the risk-equivalent value of its measured future compute saving, then $U_m<\tau\Delta C_m$ is a finite-sample sufficient condition for compute-aware dominance among the declared candidate continuations.
\end{corollary}

\begin{proof}
Immediate from Theorem~\ref{thm:paired-damage} and the definition of $b_m$. The compute-aware specialization substitutes the same risk/compute budget used in Corollary~\ref{cor:compute-dominance}.
\end{proof}

\begin{remark}[Why pairing matters]
The dense and reduced predictors are evaluated on the same examples. The variance entering the certificate is therefore
\[
\mathrm{Var}\!\left(\ell_B(f_m,Z)-\ell_B(f_0,Z)\right),
\]
not the sum of two unrelated risk-estimation variances. When the two predictors make highly correlated errors, this paired variance can be much smaller. The structural-decision problem should exploit that cancellation.
\end{remark}

\subsection{A distribution-free observability barrier}
The preceding theorem is valid but can be operationally useless when the structural effect is tiny relative to the global loss range. This is not merely a loose constant in one concentration inequality. There is a rare-event obstruction for \emph{every} distribution-free certificate that knows only a global range.

\begin{proposition}[Rare-event lower barrier for bounded-risk certification]\label{prop:range-barrier}
Let $U_n(Z_1,\ldots,Z_n)$ be any deterministic one-sided upper confidence bound satisfying
\[
\Pr_P\!\left\{\mathbb E_P Z\le U_n(Z_1,\ldots,Z_n)\right\}\ge1-\delta
\]
for every distribution $P$ supported on $[-R,R]$. Then on the all-zero sample path,
\[
\boxed{
U_n(0,\ldots,0)
\ge
R\left(1-\delta^{1/n}\right).
}
\]
Consequently, to make certification at budget $b<R$ even possible on an all-zero probe, it is necessary that
\[
\boxed{
n\ge
\frac{\log\delta}{\log(1-b/R)}
\sim
\frac{R}{b}\log\frac1\delta
\qquad (b/R\downarrow0).
}
\]
\end{proposition}

\begin{proof}
Suppose $U_n(0,\ldots,0)<R(1-\delta^{1/n})$. Choose
\[
\frac{U_n(0,\ldots,0)}R<p<1-\delta^{1/n}
\]
and let $P_p$ put mass $p$ at $R$ and $1-p$ at $0$. Then $\mathbb E_{P_p}Z=pR>U_n(0,\ldots,0)$, while the all-zero sample occurs with probability $(1-p)^n>\delta$. Hence the claimed $(1-\delta)$ coverage fails. The sample-size inequality is obtained by requiring $R(1-\delta^{1/n})\le b$ and solving for $n$; the asymptotic expression follows from $\log(1-x)\sim-x$.
\end{proof}

\begin{remark}[Interpretation]
A probe containing no visible damage is still compatible with a rare event of probability of order $\log(1/\delta)/n$ carrying loss difference of order $R$. Thus no amount of variance cancellation on the observed sample removes the $R/n$ obstruction without additional information. The correct response is not threshold tuning; it is to localize the plausible structural-damage scale using the geometry of the contraction.
\end{remark}

\subsection{Geometry-localized observability}
For clipped square loss, the master prediction-gap theorem supplies exactly the additional information needed to localize the paired loss difference.

Define
\[
\ell_B(a,y):=\min\{(a-y)^2,B\}.
\]

\begin{lemma}[Prediction gap controls paired-loss second moment]\label{lem:loss-gap-moment}
For every $a,b,y\in\mathbb R$,
\[
|\ell_B(a,y)-\ell_B(b,y)|
\le
2\sqrt B\,|a-b|.
\]
Consequently, if predictors $f_m,f_0$ satisfy
\[
\norm{f_m-f_0}_{L^2(P_X)}\le d_m,
\]
then their paired clipped-loss difference obeys
\[
\boxed{
\mathbb E Z_m^2
\le
M_{2,m}:=\min\{B^2,4Bd_m^2\}.
}
\]
\end{lemma}

\begin{proof}
The scalar function $r\mapsto\min\{r^2,B\}$ is globally $2\sqrt B$-Lipschitz: on $|r|<\sqrt B$ its derivative has magnitude at most $2\sqrt B$, and outside that interval it is constant. Apply this with $r=a-y$ and $r=b-y$. Squaring and averaging gives $\mathbb E Z_m^2\le4B\mathbb E(f_m-f_0)^2$. The trivial bound $|Z_m|\le B$ gives the minimum with $B^2$.
\end{proof}

A second-moment bound alone does not imply a smaller almost-sure range, but it can be converted into one after deterministic truncation while paying an explicit tail remainder.

For $c\in(0,B]$ write
\[
T_c(z):=\max\{-c,\min\{z,c\}\},
\qquad
W_{i,m}^{(c)}:=T_c(Z_{i,m}).
\]

\begin{theorem}[Truncated second-moment localized structural-damage bound]\label{thm:localized-damage}
Suppose that, before observing the probe, deterministic bounds $M_{2,m}\ge\mathbb E Z_m^2$ and truncation levels $c_m\in(0,B]$ are fixed for $m=1,\ldots,M$. Let $\widehat W_m$ and $\widehat v_{m,c}$ denote the empirical mean and unbiased variance of $W_{i,m}^{(c_m)}$. Then with probability at least $1-\delta$, simultaneously for all candidates,
\[
\boxed{
\begin{aligned}
D_m\le
\widehat W_m
&+\sqrt{\frac{2\widehat v_{m,c}\log(2M/\delta)}{n}}
+\frac{14c_m\log(2M/\delta)}{3(n-1)}\\
&+\mathbf 1\{c_m<B\}\frac{M_{2,m}}{4c_m}.
\end{aligned}}
\]
If $c_m=B$, the result reduces to the ordinary paired empirical-Bernstein certificate. If $c_m<B$, the deterministic truncation scale may be much smaller than the global loss range.
\end{theorem}

\begin{proof}
For $c<B$ and every real $z$,
\[
z\le T_c(z)+\frac{z^2}{4c}.
\]
Indeed the inequality is immediate for $z\le c$ because truncation of the negative tail only increases $z$. For $z>c$ it is equivalent to
\[
z-c\le\frac{z^2}{4c},
\]
which follows from $(z-2c)^2\ge0$. Therefore
\[
D_m=\mathbb E Z_m
\le
\mathbb E W_m^{(c_m)}+\frac{M_{2,m}}{4c_m}.
\]
Now $W_m^{(c_m)}\in[-c_m,c_m]$, so the same empirical-Bernstein argument as in Theorem~\ref{thm:paired-damage}, followed by a union bound over the fixed candidate family, gives the first three terms. When $c_m=B$, no truncation occurs and no moment remainder is needed.
\end{proof}

The deterministic range term and moment remainder can be balanced without looking at the probe. Ignoring the empirical-variance term, the minimizer is
\[
\boxed{
c_m^\star
=
\min\left\{
B,
\sqrt{\frac{3(n-1)M_{2,m}}{56\log(2M/\delta)}}
\right\}.}
\]
When the cap does not bind, the two deterministic terms collapse to
\[
\sqrt{\frac{14M_{2,m}\log(2M/\delta)}{3(n-1)}}.
\]
Thus a small theory-derived second moment converts the global $B/n$ obstruction into a local scale of order $\sqrt{M_{2,m}/n}$.

\begin{corollary}[Geometry-localized finite-sample certificate]\label{cor:geometry-localized}
Let $\mathfrak D_{t,H}(S_m)$ denote the terminal $L^2(P_X)$ prediction-discrepancy upper bound supplied by Theorem~\ref{thm:master-safe} for structural action $S_m$; explicitly it is the bracket in that theorem multiplied by $B_J$. Set
\[
M_{2,m}^{\rm geom}
:=
\min\left\{B^2,4B\mathfrak D_{t,H}(S_m)^2\right\}.
\]
Then Theorem~\ref{thm:localized-damage} with $M_{2,m}=M_{2,m}^{\rm geom}$ gives a valid finite-sample upper bound on the population clipped-loss structural damage. A candidate is certified whenever this localized UCB is below its predeclared risk/compute budget.
\end{corollary}

\begin{proof}
The master theorem bounds the terminal prediction discrepancy in the prediction Hilbert norm. Apply Lemma~\ref{lem:loss-gap-moment}, then Theorem~\ref{thm:localized-damage}.
\end{proof}

\begin{remark}[What changed]
The generic probe certificate asks the data to rule out every bounded rare-event distribution. That is unnecessarily pessimistic after the geometry has already certified that dense and reduced predictors must stay close in $L^2$. The localized theorem uses the population geometry to constrain the second moment and uses the independent probe only to resolve the remaining signed risk difference. This creates a direct bridge
\[
\boxed{
\text{spectral safety}
\;\longrightarrow\;
\text{prediction-gap control}
\;\longrightarrow\;
\text{localized observability}
\;\longrightarrow\;
\text{compute decision}.}
\]
\end{remark}

\begin{remark}[Why a good probe forecaster can still be useless]
A model may predict $\widehat D_m$ accurately across checkpoints while $\widehat D_m$ itself is too noisy to resolve $D_m$. The relevant decision object is therefore an upper confidence bound, not point-prediction $R^2$. This distinction is exactly what the small development audit exposed.
\end{remark}

\begin{remark}[Adaptive use of the probe]
All candidate families, moment bounds, and truncation levels above are fixed independently of the probe on which the bound is evaluated. If the same probe is used to alter future network parameters or repeatedly construct new candidates, that independence is lost. The fresh-probe oracle in the next section is one clean remedy; cross-fitting, stability, confidence-sequence, or reusable-holdout arguments are alternatives.
\end{remark}

\section{From geometric motion to a decision variable}
The first proxy implementation used total principal-angle speed of a fixed-rank hidden-feature subspace as a hard safety gate. The controlled pilot shows why that is too conservative: weak directions can rotate substantially even when virtually all current target accessibility is already saturated. The theory itself predicts this failure.

Let $P_t$ be a differentiable representation projector, let
\[
g_t=P_tf^\star,\qquad r_t=(I-P_t)f^\star,
\qquad
q_t=\frac{\norm{g_t}^2}{\norm{f^\star}^2},
\qquad
\psi_t=\arcsin\sqrt{q_t}.
\]
The feature-value identity gives
\[
\dot q_t
=
\frac{2\langle r_t,\dot P_tg_t\rangle}{\norm{f^\star}^2},
\qquad
|\dot\psi_t|\le \norm{\dot P_t}_{\op}.
\]
The inequality is a safety envelope, not an equality and not a good generic ranking statistic.

\begin{proposition}[Total subspace speed can be pure nuisance motion]\label{prop:nuisance-motion}
Fix any $q\in(0,1]$ and any $M>0$. In a Hilbert space of sufficiently large finite dimension there exists a continuously differentiable fixed-rank projector path $P_t$ and a nonzero target $f^\star$ such that
\[
q_t\equiv q,
\qquad
\dot q_t\equiv0,
\qquad
\sup_t\norm{\dot P_t}_{\op}\ge M.
\]
Hence arbitrarily rapid Grassmannian motion can carry exactly zero representation value for the target.
\end{proposition}

\begin{proof}
Choose orthogonal unit vectors $e_1,e_2$ and write $f^\star=\sqrt q\,e_1+\sqrt{1-q}\,e_2$. Let $P_t$ contain $e_1$, exclude $e_2$, and contain an additional nuisance line $v_t$ rotating with angular velocity $M$ inside a two-dimensional subspace orthogonal to $\mathrm{span}\{e_1,e_2\}$. Then $P_tf^\star=\sqrt q\,e_1$ for all $t$, so $q_t=q$ and $\dot q_t=0$, whereas the rotating nuisance line gives $\norm{\dot P_t}_{\op}=M$.
\end{proof}

This proposition explains why individual unresolved-bulk eigenvectors must not control contraction. Their rotations can be statistically unstable and target-irrelevant. The correct hierarchy is:
\[
\boxed{
\text{resolved cluster motion}
\;\longrightarrow\;
\text{target-weighted accessibility motion}
\;\longrightarrow\;
\text{finite-horizon value}.}
\]

\begin{proposition}[Architecture-free saturation certificate]\label{prop:saturation}
For every future representation projector $P_{t+H}$, regardless of its path,
\[
V^{\rm rep}_{t,H}
=\norm{f^\star}^2(q_{t+H}-q_t)
\le
(1-q_t)\norm{f^\star}^2.
\]
Thus if $q_t\ge1-\varepsilon$, no future representation path can create more than $\varepsilon\norm{f^\star}^2$ of additional population span value.
\end{proposition}

\begin{proof}
Every orthogonal projection satisfies $q_{t+H}\le1$.
\end{proof}

The sharper finite-horizon certificate uses a valid path-length upper bound $L_{t,H}$:
\[
V^{\rm rep}_{t,H}
\le
\norm{f^\star}^2
\left[
\sin^2\!\left(
\min\left\{\frac\pi2,\arcsin\sqrt{q_t}+L_{t,H}\right\}
\right)-q_t
\right].
\]
A recent observed target-angle or Grassmannian speed may be inserted as a forecast, but then the quantity is a \emph{proxy}, not a certificate. The no-universal-finite-jet result from the endogenous spectral theory is relevant here: current local derivatives alone cannot provide an architecture-free exact bound on all future accelerations of the learned geometry.

\section{Algorithm derived from the theory}
The theoretical controller has five gates. At sparse checkpoints, estimate a stable target-weighted prediction/representation state and compare a finite family of structural actions. The gates are deliberately ordered: a candidate that fails an early, cheap gate never reaches the more expensive counterfactual checks.

\paragraph{Gate A: current finite-horizon mode value.}
For an ideal fixed-geometry mode, use
\[
\widehat V_j(H)=\frac12\widehat a_j^2
\left(1-e^{-2\widehat\mu_jH}\right).
\]
For a structural candidate (e.g. neuron subset), measure its current target-accessibility loss and matched frozen-readout distortion. Eigenvalue magnitude by itself is never sufficient.

\paragraph{Gate B: future discovery value.}
Do not track arbitrary individual bulk eigenvectors. First define a \emph{resolved} spectral cluster by a predeclared null/edge rule or another stability-certified cluster criterion. Measure its target accessibility and target-angle progress. Use, in decreasing order of rigor:
\begin{enumerate}[label=(\alph*)]
\item an architecture-specific certified finite-horizon path-length/value bound;
\item the architecture-free saturation bound of Proposition~\ref{prop:saturation};
\item a calibrated prospective value forecaster;
\item only as a diagnostic, recent target-angle or resolved-cluster Grassmannian speed.
\end{enumerate}
A large rotation confined to unresolved nuisance directions is not a reason to keep the full network dense.

\paragraph{Gate C: structural realizability.}
Require a small structural approximation term $\delta$ in Theorem~\ref{thm:structural-transfer}. For the final hidden layer this can be implemented by rank-revealing or target-aware neuron subset selection followed by the same readout solve in both the reduced branch and the full-width control. For an internal nonlinear layer, no spectral action is declared realizable until a channel/head/rank deletion is shown to approximate the desired kept parameter cluster.

\paragraph{Gate D: finite-sample observability.}
A point estimate of probe damage is not enough. For every candidate continuation $m$, evaluate the reduced and matched dense predictors on the \emph{same} independent probe examples, form the paired loss differences, and compute the upper confidence bound from Theorem~\ref{thm:paired-damage}:
\[
U_m
=
\widehat D_m
+
\sqrt{\frac{2\widehat v_m\log(2M/\delta)}{n}}
+
\frac{14B\log(2M/\delta)}{3(n-1)}.
\]
A candidate survives only if its structural damage is statistically resolved tightly enough that $U_m$ lies inside its predeclared risk budget. If the confidence radius is too wide, the correct action is \emph{do not contract yet} or acquire a larger/variance-reduced probe; it is not to trust the sign of $\widehat D_m$.

\paragraph{Gate E: compute value.}
Let $\Delta C_m$ be the measured future compute saving of candidate $m$, including all controller, probe, readout, and rebuilding overhead. For risk-equivalent compute price $\tau$, contract only if
\[
\boxed{U_m<\tau\Delta C_m}
\]
when a paired finite-sample certificate is available. More generally, use the strongest valid upper bound on structural risk price and require it to be below the compute value. This is the empirical counterpart of Corollary~\ref{cor:compute-dominance}.

\begin{center}
\fbox{\parbox{0.92\linewidth}{
\textbf{Certified Target-Value Spectral Rank Annealing (TV-SRA).}
Train the dense network during the discovery phase. At sparse checkpoints, identify statistically resolved target-bearing geometry and construct a small finite family of structurally realizable contractions. Bound the remaining future target value, then bound the paired finite-sample structural damage of each surviving action. Contract only when a candidate's structural-damage upper confidence bound is below its measured risk-equivalent compute saving. Rebuild only the affected optimizer state and continue the same base optimizer. If no representation block retains positive value-per-compute, freeze the backbone and solve the readout.
}}
\end{center}

The current repository implementation intentionally contracts only the final hidden layer. This is the cleanest place where a coordinate-level neuron deletion produces a genuine parameter/FLOP reduction and where the representation/prediction spectrum has an exact linear-readout interpretation. Extension to internal nonlinear layers requires an explicit structural-transfer bound rather than an analogy.

The finite-horizon interpretation is consistent with the adjoint control formulation of the endogenous spectral-geometry theory. If a reduced geometry state obeys $q_{s+1}=F_s(q_s,u_s)$, an irreversible contraction is a discrete control action whose value must be judged through terminal risk, not through an instantaneous spectral descent direction. TV-SRA is the coarse structural-control specialization: it leaves the base optimizer untouched until the future value of a capacity sector is smaller than the compute that sector costs.

\section{Adaptive probe safety for repeated contractions}
A one-shot independent probe is not automatically valid after the first contraction, because the contraction changes all future checkpoints. The cleanest theorem uses a fresh probe at every irreversible decision.

\begin{theorem}[Fresh-probe adaptive contraction oracle]\label{thm:fresh-probe}
Consider decision times $k=1,\ldots,K$. Conditional on all data, models, and contraction decisions before time $k$, suppose the controller constructs a finite candidate family $\{f_{k,m}:1\le m\le M_k\}$ without using the fresh probe $\mathcal P_k$ of size $n_k$. Let the probe loss be clipped to $[0,B]$, let $c_{k,m}$ be a deterministic compute penalty measurable before seeing $\mathcal P_k$, and select
\[
\widehat m_k\in\arg\min_m\{\widehat R_{\mathcal P_k}(f_{k,m})+c_{k,m}\}.
\]
For any numbers $\delta_k>0$ with $\sum_{k=1}^K\delta_k\le\delta$, with probability at least $1-\delta$ simultaneously for every adaptive decision time,
\[
\boxed{
R(f_{k,\widehat m_k})+c_{k,\widehat m_k}
\le
\min_{1\le m\le M_k}\{R(f_{k,m})+c_{k,m}\}
+2B\sqrt{\frac{\log(2M_k/\delta_k)}{2n_k}}.}
\]
\end{theorem}

\begin{proof}
Condition on the entire history before decision $k$. The candidate predictors and compute penalties are then fixed, while $\mathcal P_k$ is independent. Hoeffding's inequality and a union bound over $M_k$ candidates imply that with conditional probability at least $1-\delta_k$ all candidate probe risks deviate from population risk by at most
\[
B\sqrt{\frac{\log(2M_k/\delta_k)}{2n_k}}.
\]
On this event, empirical minimization plus two uniform deviations gives the displayed oracle inequality. A union bound over adaptive decision times completes the proof.
\end{proof}

This theorem gives a rigorous repeated-decision protocol. Reusing the same probe indefinitely requires an additional stability, cross-fitting, reusable-holdout, or martingale argument and is therefore treated as a practical proxy rather than a certificate.

\section{Empirical protocol and falsification criteria}
The decisive comparison is
\[
\text{same initialization + same data order + same base optimizer}
\]
for a dense model and the contraction controller. The primary plots are test metric versus total measured compute and active parameters versus training time. Every SVD, probe forward pass, Jacobian call, checkpoint, readout solve, confidence-bound evaluation, and model-rebuild cost is charged to the proposed method.

The minimum empirical ladder is:
\begin{enumerate}
\item \textbf{Exact fixed-geometry audit.} Numerically verify Theorem~\ref{thm:fixed-value} and the compute-to-accuracy consequence to machine precision.
\item \textbf{Solvable quadratic teacher--student.} Verify the exact discovery-delay law: blind contraction before target resolution starts with vanishing overlap and pays a dimension-dependent delay, while post-BBP contraction retains an order-one aligned direction.
\item \textbf{Finite-sample observability audit.} On a model where population structural damage can be approximated with a very large independent evaluation sample, check nominal coverage and tightness of the paired empirical-Bernstein bound as probe size grows. This must precede use of the UCB as a controller gate.
\item \textbf{Small real MLPs.} Digits and breast cancer are development-only after the first two controller iterations. Compare dense AdamW, optimizer-reset, dense readout-refit, the falsified total-speed SRA proxy, and TV-SRA.
\item \textbf{Image scale.} ResNet/CIFAR and then ImageNet. Contract channels, not unstructured weights, so FLOP savings are real. EPI and Early-Bird are mandatory onset baselines.
\item \textbf{Transformer scale.} Fine-tuning and continued pretraining. Candidate structural units are MLP intermediate rank, attention heads, or low-rank factors. Compare with dynamic sparse/rank and modern adaptive-PEFT baselines.
\end{enumerate}

A result is considered negative if any of the following holds after matched compute: the controller does not reduce wall-clock/FLOPs, test performance falls outside a predeclared tolerance, probe overhead consumes the saved compute, structural-damage confidence intervals are too wide to trigger useful contractions, or the onset rule fails to transfer across configurations.

\section{Position relative to existing methods}
The novelty cannot be ``pruning during training,'' ``adaptive rank,'' or simply ``when to prune.'' Early-Bird tickets and, most directly, the Early Pruning Indicator (EPI) of Shen et al. already detect an early point at which a dominant structural-pruning architecture stabilizes. RigL/SRigL, EigenDamage, spectral pruning, dynamic rank methods, AdaLoRA, GaLore, FARSS, and SPARCS occupy adjacent spaces.

The distinctive scientific claim pursued here is narrower:
\begin{quote}
A target-weighted prediction-space theory determines the finite-horizon value of capacity, proves a mechanism by which unresolved target directions become identifiable only after a training-generated phase transition, supplies an explicit error budget for turning that prediction-space decision into structural contraction, and requires the resulting structural damage to be statistically observable before compute is traded for capacity.
\end{quote}
The empirical distinction from EPI must be tested directly: architecture/mask stability versus target-value/geometry onset, under identical structural pruning actions and charged compute.

\section{Claim ledger}
\begin{center}
\begin{tabular}{p{0.42\linewidth}p{0.16\linewidth}p{0.32\linewidth}}
\toprule
Claim & Status & Scope \\
\midrule
Modal gradient energy $\sum\mu_ja_j^2$ & Exact & Any differentiable square-loss predictor \\
Fixed-geometry deletion cost & Exact & Frozen Jacobian / linearized dynamics \\
Compute-to-target bound & Exact & Fixed geometry above deletion floor \\
Value-per-compute oracle & Exact & Fixed geometry, separable mode costs \\
Finite-sample drop criterion & Exact & Gaussian sequence model \\
Total subspace speed can be pure nuisance motion & Exact & Projector geometry \\
Accessibility saturation bound & Exact & Arbitrary future representation path \\
Curvature-to-spectral path bound & Exact conditional & Bounded Jacobian and second derivative \\
Structural-flow stability & Exact conditional & Lipschitz tube + omitted-gradient envelope \\
Spectral persistence bound & Exact conditional & Isolated cluster + operator-velocity/gap bounds \\
Master structural risk bound & Exact conditional & Geometry motion + structural mismatch + jump \\
Paired structural-damage UCB & Exact finite-sample & Fixed candidate family + independent probe \\
Exact neuron-basis contraction & Exact & Current final-hidden linear-readout span \\
Dynamic target compressibility after BBP & Exact in model & Quadratic teacher--student phase theorem \\
Rank-one discovery-delay law & Exact in model & Post-contraction target-angle dynamics \\
Fresh-probe adaptive oracle & Exact & Finite candidate grids + fresh independent probes \\
Universal internal-layer contraction & Open & Requires architecture-specific structural bridge \\
Compute-matched speedup & Empirical question & Must be established, not assumed \\
\bottomrule
\end{tabular}
\end{center}

\section{What remains mathematically hard}
Four problems separate this draft from a definitive theory. First, a cheap observable upper bound on future \emph{target-weighted} geometry motion is architecture-dependent; total Grassmann speed is universally safe but can be arbitrarily loose by Proposition~\ref{prop:nuisance-motion}. Second, converting a prediction-space cluster into a hardware-friendly internal-layer contraction without large structural gap $\delta$ requires architecture-specific results. Third, the paired finite-sample certificate can be too conservative when the probe is small or the loss bound is loose; variance-adaptive or problem-specific concentration may be needed. Fourth, the fresh-probe theorem gives a clean repeated-decision route, but repeatedly spending fresh labels can be expensive; a reusable-probe or cross-fitting guarantee with low sample overhead remains open.

These are not reasons to avoid the algorithm. They are the exact places where the theory says empirical shortcuts must be labeled as proxies rather than certificates.

\begin{thebibliography}{99}
\bibitem{BignozziESG}
E. Bignozzi. \emph{Endogenous Spectral Geometry of Deep Learning}. Working research monograph, 2026.

\bibitem{BignozziFVL}
E. Bignozzi. \emph{The Value of Feature Learning: Geometry-Aware Forecasts for Compute-Efficient Neural Training}. Working paper, 2026.

\bibitem{DavisKahan}
C. Davis and W. M. Kahan. The rotation of eigenvectors by a perturbation. III. \emph{SIAM Journal on Numerical Analysis}, 1970.

\bibitem{Kato}
T. Kato. \emph{Perturbation Theory for Linear Operators}. Springer, 1995.

\bibitem{EarlyBird}
H. You et al. Drawing Early-Bird Tickets: Toward More Efficient Training of Deep Networks. ICLR, 2020.

\bibitem{EPI}
M. Shen, P. Molchanov, H. Yin, and J. M. Alvarez. When To Prune? A Policy Towards Early Structural Pruning. CVPR, 2022.

\bibitem{RigL}
U. Evci, T. Gale, J. Menick, P. S. Castro, and E. Elsen. Rigging the Lottery: Making All Tickets Winners. ICML, 2020.

\bibitem{SRigL}
M. Lasby, A. Golubeva, U. Evci, M. Nica, and Y. Ioannou. Dynamic Sparse Training with Structured Sparsity. ICLR, 2024.

\bibitem{EigenDamage}
C. Wang, R. Grosse, S. Fidler, and G. Zhang. EigenDamage: Structured Pruning in the Kronecker-Factored Eigenbasis. ICML, 2019.

\bibitem{SpectralPruning}
L. Buffoni et al. Spectral pruning of fully connected layers. \emph{Scientific Reports}, 2022.

\bibitem{AdaLoRA}
Q. Zhang et al. AdaLoRA: Adaptive Budget Allocation for Parameter-Efficient Fine-Tuning. 2023.

\bibitem{GaLore}
J. Zhao et al. GaLore: Memory-Efficient LLM Training by Gradient Low-Rank Projection. ICML, 2024.

\bibitem{DynamicRank}
H. Shin, A. Jung, S. Lee, and S. Hong. Dynamic Rank Adjustment for Accurate and Efficient Neural Network Training. 2025.

\bibitem{FARSS}
Y. Huang et al. FARSS: Fisher-Optimized Adaptive Low-Rank and Singular-Vector Selection for Knowledge-Preserving Fine-Tuning. Findings of ACL, 2026.

\bibitem{SPARCS}
G. Peri et al. SPectral ARchiteCture Search for neural network models. \emph{npj Artificial Intelligence}, 2025.
\end{thebibliography}

\end{document}

\end{document}